% !LW recipe = latexmk (xelatex)
% report.tex -- Báo cáo Assignment 1 (Parallel Computing, HCMUT).
%
%   Bố cục theo template báo cáo HCMUT: https://github.com/truongnh28/hcmut-report-template
%   (header có logo, footer có số trang, trang bìa, mục "Member list & Workload", "References").
%   Template gốc dùng pdfLaTeX + vntex; báo cáo này dùng XeLaTeX để in được chữ Unicode trong các
%   bảng do scripts/plot.py sinh ra và trong mã nguồn, với phông Latin Modern (cùng họ phông
%   Computer Modern của template).
%
%   Biên dịch:  make report        (= cd report && latexmk -xelatex report.tex)
%   Cần XeLaTeX và TeX Live đầy đủ; phông chữ lấy từ TeX Live, không phụ thuộc phông của hệ điều hành.
%
% ĐÂY LÀ BỘ KHUNG. Mọi hình, bảng số liệu và chú thích của chúng do scripts/plot.py sinh ra từ
% results/<máy>/ (make figures) và được nạp qua \PCfigure, \PCtable; báo cáo không chứa con số
% nào chép tay. Chỗ nào còn phải viết được đánh dấu "% TODO(report):" trong mã nguồn và hiện
% thành ghi chú màu xám trong PDF (tắt bằng \showtodofalse ở dưới).
\documentclass[a4paper]{article}

% =================================================================================================
% THÔNG TIN CỦA NHÓM -- sửa ở đây
% =================================================================================================
\newcommand{\ReportTitle}{Mutex Lock, Sense-Reversing Barrier\\ và Concurrent Counter}
\newcommand{\ReportSubtitle}{Cài đặt, kiểm tra tính đúng đắn và nghiên cứu hiệu năng}
\newcommand{\CourseName}{Tính toán song song (Parallel Computing)}
\newcommand{\AssignmentName}{Assignment 1}
\newcommand{\Lecturers}{PGS.TS. Trần Văn Hoài, TS. Diệp Thanh Đăng}
\newcommand{\GroupName}{Nhóm \ldots}                       % TODO(report): tên hoặc số nhóm
% TODO(report): họ tên và mã số học viên của từng thành viên, mỗi người một dòng.
% Bảng "Students" của trang bìa, như template: dòng đầu có chữ "Students:", các dòng sau để trống
% hai cột đầu. Mỗi dòng: <Họ và tên> - <MSHV>.
\newcommand{\GroupMembers}{%
  & Students: & Nguyễn Hữu Trưởng - 2470573.\\
  & & Nguyễn Nhật Thanh - 2570316.\\
  & & Võ Thị Bích Phượng - 2470570.\\
}
\newcommand{\CourseTitleEN}{PARALLEL COMPUTING}          % dòng tên môn trên trang bìa
\newcommand{\AdvisorLine}{Assoc. Prof. Trần Văn Hoài\\ & & PhD. Diệp Thanh Đăng}
\newcommand{\AcademicYear}{2026 - 2027}
\newcommand{\ReportDate}{HO CHI MINH CITY, OCTOBER 2026}

% Các máy đã đo: tên thư mục trong results/, cách nhau bằng dấu phẩy, KHÔNG có dấu cách.
% Thêm một máy = thêm tên của nó vào dòng này (và một dòng \PCsetowner bên dưới).
\providecommand{\Hosts}{apple-m5,C4M7Y670NT,LAP1184}
% Thư mục chứa hình và bảng do scripts/plot.py sinh ra.
\providecommand{\FigRoot}{figures}
% =================================================================================================

\usepackage{a4wide}
\usepackage{amsmath,amssymb}
\usepackage{fontspec}
\usepackage{unicode-math}
% Latin Modern: cùng họ phông Computer Modern của template, đủ chữ tiếng Việt.
\setmainfont{lmroman10}[Extension=.otf, UprightFont=*-regular, BoldFont=*-bold,
  ItalicFont=*-italic, BoldItalicFont=*-bolditalic]
\setsansfont{lmsans10}[Extension=.otf, UprightFont=*-regular, BoldFont=*-bold,
  ItalicFont=*-oblique, BoldItalicFont=*-boldoblique]
\setmonofont{lmmono10}[Extension=.otf, UprightFont=*-regular, BoldFont=lmmonolt10-bold,
  ItalicFont=*-italic, BoldItalicFont=lmmonolt10-boldoblique]
\setmathfont{latinmodern-math.otf}
\usepackage[vietnamese]{babel}
% Như template: mục lục và tài liệu tham khảo mang tên tiếng Anh.
\addto\captionsvietnamese{\renewcommand*\contentsname{Contents}\renewcommand*\refname{References}}

\usepackage{graphicx}
\usepackage{booktabs}
\usepackage{array}
\usepackage{tabularx}
\usepackage{multirow}
\usepackage{float}
\usepackage[font=small,labelfont=bf]{caption}
\usepackage{xcolor}
\usepackage{enumitem}
\usepackage{etoolbox}
\usepackage{listings}
\usepackage{lastpage}
\usepackage{fancyhdr}
\usepackage{hyperref}
\hypersetup{urlcolor=blue,linkcolor=black,citecolor=black,colorlinks=true}

% Hộp trả lời cho từng câu hỏi nghiên cứu (RQ) ở mục Kết quả.
\usepackage{tcolorbox}
\definecolor{rqblue}{HTML}{E3ECF5}
\newtcolorbox{rqanswer}{colback=rqblue,colframe=rqblue!60!black,boxrule=0.4pt,arc=1pt,
  left=5pt,right=5pt,top=3pt,bottom=3pt}

% ---- header và footer của template ----
\setlength{\headheight}{40pt}
\pagestyle{fancy}
\fancyhead{}
\fancyhead[L]{%
 \begin{tabular}{rl}
    \begin{picture}(25,15)(0,0)
    \put(0,-8){\includegraphics[width=8mm, height=8mm]{hcmut.png}}
   \end{picture}&
	\begin{tabular}{l}
		\textbf{\ttfamily University of Technology, Ho Chi Minh City}\\
		\textbf{\ttfamily Faculty of Computer Science and Engineering}
	\end{tabular}
 \end{tabular}}
\fancyhead[R]{}
\fancyfoot{}
\fancyfoot[L]{\scriptsize\ttfamily Assignment for Parallel Computing - Academic year \AcademicYear}
\fancyfoot[R]{\scriptsize\ttfamily Page {\thepage}/\pageref{LastPage}}
\renewcommand{\headrulewidth}{0.3pt}
\renewcommand{\footrulewidth}{0.3pt}
\setcounter{secnumdepth}{4}
\setcounter{tocdepth}{2}

\setlength{\parskip}{0.4em}
\setlength{\emergencystretch}{2em}

% ---- mã nguồn trong báo cáo ----
% Style "myStyle" của template (màu nền, màu từ khoá, chú thích, chuỗi).
\definecolor{codegreen}{rgb}{0,0.6,0}
\definecolor{codegray}{rgb}{0.5,0.5,0.5}
\definecolor{codepurple}{rgb}{0.58,0,0.82}
\definecolor{backcolour}{rgb}{0.95,0.95,0.92}
\lstset{language=C++, backgroundcolor=\color{backcolour}, commentstyle=\color{codegreen},
  keywordstyle=\color{magenta}, numberstyle=\tiny\color{codegray}, stringstyle=\color{codepurple},
  basicstyle=\ttfamily\footnotesize, columns=fullflexible, keepspaces=true, numbers=left,
  numbersep=5pt, showspaces=false, showstringspaces=false, showtabs=false, tabsize=2,
  breakatwhitespace=false, breaklines=true, aboveskip=0.8em, belowskip=0.6em, extendedchars=true}
% listings dưới XeLaTeX xếp sai thứ tự các ký tự ngoài ASCII (ví dụ chữ có dấu trong chú thích).
% Khai báo các chữ cái tiếng Việt là "chữ" của listings để chúng được in đúng chỗ.
\makeatletter
\lst@InputCatcodes
\def\lst@DefEC{%
  \lst@CCECUse \lst@ProcessLetter
  ^^80^^81^^82^^83^^84^^85^^86^^87^^88^^89^^8a^^8b^^8c^^8d^^8e^^8f%
  ^^90^^91^^92^^93^^94^^95^^96^^97^^98^^99^^9a^^9b^^9c^^9d^^9e^^9f%
  ^^a0^^a1^^a2^^a3^^a4^^a5^^a6^^a7^^a8^^a9^^aa^^ab^^ac^^ad^^ae^^af%
  ^^b0^^b1^^b2^^b3^^b4^^b5^^b6^^b7^^b8^^b9^^ba^^bb^^bc^^bd^^be^^bf%
  ^^c0^^c1^^c2^^c3^^c4^^c5^^c6^^c7^^c8^^c9^^ca^^cb^^cc^^cd^^ce^^cf%
  ^^d0^^d1^^d2^^d3^^d4^^d5^^d6^^d7^^d8^^d9^^da^^db^^dc^^dd^^de^^df%
  ^^e0^^e1^^e2^^e3^^e4^^e5^^e6^^e7^^e8^^e9^^ea^^eb^^ec^^ed^^ee^^ef%
  ^^f0^^f1^^f2^^f3^^f4^^f5^^f6^^f7^^f8^^f9^^fa^^fb^^fc^^fd^^fe^^ff%
  ^^^^0102^^^^0103^^^^0110^^^^0111^^^^0128^^^^0129^^^^0168^^^^0169^^^^01a0^^^^01a1^^^^01af%
  ^^^^01b0^^^^1ea0^^^^1ea1^^^^1ea2^^^^1ea3^^^^1ea4^^^^1ea5^^^^1ea6^^^^1ea7^^^^1ea8^^^^1ea9%
  ^^^^1eaa^^^^1eab^^^^1eac^^^^1ead^^^^1eae^^^^1eaf^^^^1eb0^^^^1eb1^^^^1eb2^^^^1eb3^^^^1eb4%
  ^^^^1eb5^^^^1eb6^^^^1eb7^^^^1eb8^^^^1eb9^^^^1eba^^^^1ebb^^^^1ebc^^^^1ebd^^^^1ebe^^^^1ebf%
  ^^^^1ec0^^^^1ec1^^^^1ec2^^^^1ec3^^^^1ec4^^^^1ec5^^^^1ec6^^^^1ec7^^^^1ec8^^^^1ec9^^^^1eca%
  ^^^^1ecb^^^^1ecc^^^^1ecd^^^^1ece^^^^1ecf^^^^1ed0^^^^1ed1^^^^1ed2^^^^1ed3^^^^1ed4^^^^1ed5%
  ^^^^1ed6^^^^1ed7^^^^1ed8^^^^1ed9^^^^1eda^^^^1edb^^^^1edc^^^^1edd^^^^1ede^^^^1edf^^^^1ee0%
  ^^^^1ee1^^^^1ee2^^^^1ee3^^^^1ee4^^^^1ee5^^^^1ee6^^^^1ee7^^^^1ee8^^^^1ee9^^^^1eea^^^^1eeb%
  ^^^^1eec^^^^1eed^^^^1eee^^^^1eef^^^^1ef0^^^^1ef1^^^^1ef2^^^^1ef3^^^^1ef4^^^^1ef5^^^^1ef6%
  ^^^^1ef7^^^^1ef8^^^^1ef9%
  ^^00}
\lst@RestoreCatcodes
\makeatother

% =================================================================================================
% Macro nối báo cáo với kết quả do scripts/plot.py sinh ra
%   \PCfigure{<máy>}{<tên hình>}   hình \FigRoot/<máy>/<tên hình>.pdf kèm chú thích do plot.py viết
%   \PCtable{<máy>}{<khoá>}        bảng <khoá> của \FigRoot/<máy>/tables.tex kèm chú thích
%   \PCval{<máy>}{<khoá>}          một giá trị (ncpu, cpu, os, compiler, ...)
%   \ForEachHost{\Macro}           gọi \Macro{<máy>} cho từng máy trong \Hosts
%   <máy> = cross: hình và bảng so sánh giữa các máy (\FigRoot/cross/)
% Hình hoặc bảng chưa có thì hiện một khung nhắc, báo cáo vẫn biên dịch được.
% Nhãn tham chiếu: fig:<máy>:<tên hình>, tab:<máy>:<khoá>.
% =================================================================================================
\newif\ifshowtodo
\showtodotrue            % \showtodofalse: ẩn mọi ghi chú "Cần viết" và khung nhắc trong PDF
\newcommand{\TODO}[1]{\ifshowtodo\par\noindent{\small\itshape\color{black!55}[Cần viết: #1]}\par\fi}
\newcommand{\PCmissing}[1]{\ifshowtodo\par\medskip\noindent\fcolorbox{black!35}{black!4}{%
  \parbox{\dimexpr\linewidth-2\fboxsep-2\fboxrule}{\small\itshape\color{black!65}#1}}\par\medskip\fi}
\newcommand{\PCname}[1]{\texttt{\detokenize{#1}}}
\makeatletter
\newcommand{\PCload}[1]{\IfFileExists{\FigRoot/#1/tables.tex}{\input{\FigRoot/#1/tables.tex}}{}}
\newcommand{\PCval}[2]{\ifcsname PC@#1@val@#2\endcsname\csname PC@#1@val@#2\endcsname\else\textit{(chưa có)}\fi}
\newcommand{\PCtable}[2]{%
  \ifcsname PC@#1@tab@#2\endcsname
    \begin{table}[H]\centering
      \caption{\csname PC@#1@tabcap@#2\endcsname}\label{tab:#1:#2}
      \csname PC@#1@tab@#2\endcsname
    \end{table}
  \else
    \PCmissing{Bảng \PCname{#2} của \PCname{#1} chưa có. Chạy \texttt{make bench} trên máy đó,
      chép \texttt{results/\detokenize{#1}/} về, rồi chạy \texttt{make figures}.}%
  \fi}
\newcommand{\PCfigure}[2]{%
  \IfFileExists{\FigRoot/#1/#2.pdf}{%
    \begin{figure}[tbp]\centering
      \includegraphics[width=0.8\linewidth]{\FigRoot/#1/#2.pdf}
      \caption{\ifcsname PC@#1@figcap@#2\endcsname\csname PC@#1@figcap@#2\endcsname
               \else Hình \PCname{#2} của \PCname{#1}.\fi}\label{fig:#1:#2}
    \end{figure}}{%
    \PCmissing{Hình \PCname{#2} của \PCname{#1} chưa có (\texttt{make figures} vẽ nó từ
      \texttt{results/\detokenize{#1}/}).}}}
% Hình và bảng không bắt buộc (chỉ có khi đã chạy phần đo tương ứng): không hiện khung nhắc.
\newcommand{\PCfigureopt}[2]{\IfFileExists{\FigRoot/#1/#2.pdf}{\PCfigure{#1}{#2}}{}}
\newcommand{\PCtableopt}[2]{\ifcsname PC@#1@tab@#2\endcsname\PCtable{#1}{#2}\fi}
\newcommand{\PCappendixtables}[1]{\ifcsname PC@#1@tab@appendix\endcsname\csname PC@#1@tab@appendix\endcsname\fi}
\newcommand{\ForEachHost}[1]{\def\PC@each{\forcsvlist{#1}}\expandafter\PC@each\expandafter{\Hosts}}
% Thành viên chạy đo trên từng máy (hiện trong bảng cấu hình máy).
% Máy do người viết báo cáo đo; hình của máy này nằm xen trong lời phân tích ở mục 6,
% máy khác chỉ xuất hiện trong hình so sánh và bảng ở mục So sánh giữa các máy.
\newcommand{\PrimaryHost}{apple-m5}
\newcommand{\PC@ifother}[2]{\ifstrequal{#2}{apple-m5}{}{#1{#2}}}
\newcommand{\ForEachOtherHost}[1]{\def\PC@eachother{\forcsvlist{\PC@ifother{#1}}}\expandafter\PC@eachother\expandafter{\Hosts}}
\newcommand{\PCsetowner}[2]{\expandafter\gdef\csname PC@owner@#1\endcsname{#2}}
\newcommand{\PCowner}[1]{\ifcsname PC@owner@#1\endcsname\csname PC@owner@#1\endcsname\else\textit{(điền tên)}\fi}
\makeatother

% TODO(report): mỗi máy một dòng, ghi tên thành viên đã chạy đo trên máy đó.
\PCsetowner{apple-m5}{\textit{[điền]}}
\PCsetowner{C4M7Y670NT}{\textit{[điền]}}
\PCsetowner{LAP1184}{\textit{[điền]}}

\ForEachHost{\PCload}
\PCload{cross}

% Hình và bảng của từng phần cho một máy; phần 6 gọi chúng cho mọi máy trong \Hosts.
% Giới hạn khoảng 20 trang: mỗi phần chỉ giữ các hình và bảng chính. Các hình còn lại vẫn được
% plot.py sinh ra trong figures/<máy>/ và nằm trong bài nộp, nhưng không đưa vào báo cáo.
\newcommand{\LockFigures}[1]{%
  \subsubsection*{Máy \PCname{#1}}
  \PCfigure{#1}{lock_slowdown}
  \PCtable{#1}{locktime}
  \PCfigure{#1}{lock_cs_length}
  \PCfigure{#1}{lock_backoff}
  \PCfigure{#1}{lock_fairness}
}
\newcommand{\BarrierFigures}[1]{%
  \subsubsection*{Máy \PCname{#1}}
  \PCfigure{#1}{barrier_rate}
  \PCtable{#1}{barrier}
}
\newcommand{\CounterFigures}[1]{%
  \subsubsection*{Máy \PCname{#1}}
  \PCfigure{#1}{counter_time}
  \PCtable{#1}{countertime}
  \PCfigure{#1}{counter_cas_failures}
}
\newcommand{\CampaignTables}[1]{\PCtable{#1}{campaign}\PCtable{#1}{lockunit}}

\begin{document}

% ---- Trang bìa (bố cục của template) --------------------------------------------------------------
\begin{titlepage}
\begin{center}
VIETNAM NATIONAL UNIVERSITY, HO CHI MINH CITY \\
UNIVERSITY OF TECHNOLOGY \\
FACULTY OF COMPUTER SCIENCE AND ENGINEERING
\end{center}

\vspace{1cm}

\begin{figure}[h!]
\begin{center}
\includegraphics[width=3cm]{hcmut.png}
\end{center}
\end{figure}

\vspace{1cm}

\begin{center}
\begin{tabular}{c}
\multicolumn{1}{c}{\textbf{{\Large \CourseTitleEN}}}\\
~~\\
\hline
\\
\multicolumn{1}{l}{\textbf{{\Large \AssignmentName}}}\\
\\
\textbf{{\Huge \begin{tabular}{@{}c@{}}\ReportTitle\end{tabular}}}\\
\\
\textbf{{\large \ReportSubtitle}}\\
\\
\hline
\end{tabular}
\end{center}

\vspace{2.5cm}

\begin{table}[h]
\begin{tabular}{rrl}
\hspace{4 cm} & Advisor: & \AdvisorLine\\
\GroupMembers
\end{tabular}
\end{table}

\vfill
\begin{center}
{\footnotesize \ReportDate}
\end{center}
\end{titlepage}

\tableofcontents
\newpage

% ====================================================================================================
% \section{Member list \& Workload}
% % TODO(report): nhóm tự điền họ tên, MSHV, phần việc, tỉ lệ đóng góp (theo bảng của template).
% % Mọi thành viên phải đọc toàn bộ báo cáo và trả lời được câu hỏi về bất kỳ phần nào.
% \begin{center}
% \begin{tabular}{|c|c|c|l|c|}
% \hline
% \textbf{No.} & \textbf{Fullname} & \textbf{Student ID} & \textbf{Problems} & \textbf{Percentage of work}\\
% \hline
% \multirow{3}{*}{1} & \multirow{3}{*}{\ldots} & \multirow{3}{*}{\ldots} & - \ldots & \multirow{3}{*}{\ldots\%}\\
%  & & & & \\
%  & & & & \\
% \hline
% \multirow{3}{*}{2} & \multirow{3}{*}{\ldots} & \multirow{3}{*}{\ldots} & - \ldots & \multirow{3}{*}{\ldots\%}\\
%  & & & & \\
%  & & & & \\
% \hline
% \multirow{3}{*}{3} & \multirow{3}{*}{\ldots} & \multirow{3}{*}{\ldots} & - \ldots & \multirow{3}{*}{\ldots\%}\\
%  & & & & \\
%  & & & & \\
% \hline
% \end{tabular}
% \end{center}
% \TODO{nhóm điền bảng phân công: họ tên, MSHV, phần việc (cài đặt, kiểm tra, đo trên máy nào, phân tích, viết báo cáo) và tỉ lệ đóng góp.}

% ====================================================================================================
\section{Giới thiệu và câu hỏi nghiên cứu}
% TODO(report): 0,5-1 trang.
%  - Mục tiêu của assignment: cài đặt mutex lock (TAS, TTAS, MCS), sense-reversing barrier và
%    concurrent counter (lock, compare-and-swap, fetch-and-add); kiểm tra tính đúng đắn; nghiên
%    cứu hiệu năng và khả năng mở rộng trên các máy shared-memory của nhóm.
%  - Câu hỏi nghiên cứu cho từng phần, viết thành câu hỏi trả lời được bằng số liệu. Ví dụ: lock
%    nào mở rộng tốt nhất ở mức tranh chấp cao; TAS sụp ở đâu; MCS có lợi từ lúc nào; backoff giúp
%    TAS hay TTAS nhiều hơn; lock tự viết đứng đâu so với lock thư viện; chi phí một barrier
%    episode tăng thế nào theo số thread; CAS chậm hơn fetch-and-add vì đâu.
%  - Phạm vi: phần 4 của đề (bản phân tán bằng MPI-3 RMA, không bắt buộc) không được thực hiện.
%  - Bố cục báo cáo.
% Mục 1. Thuật ngữ, khung giải thích (lưu lượng cache coherence, scheduler, hai loại core) và sáu câu hỏi nghiên cứu.
Nhóm cài đặt bằng C++ mutex lock (TAS, TTAS, MCS), sense-reversing barrier và concurrent counter
(lock, CAS, fetch-and-add). Nhóm kiểm tra tính đúng đắn của các cài đặt và đo khả năng mở rộng trên máy
shared-memory, so với cấu trúc của thư viện. Phần 4 của đề (MPI-3 RMA~\cite{mpi41}) không bắt buộc và nhóm không làm.

\begin{itemize}[nosep]
  \item \emph{Cache line.} Đơn vị (128 byte trên hai máy Apple, 64 byte trên \PCname{LAP1184})
    mà phần cứng giữ nhất quán giữa các core. Nhiều core cùng đọc được, nhưng ghi thì một core phải giữ
    line độc quyền và làm mất hiệu lực bản sao ở core khác.
  \item \emph{Tranh chấp.} Nhiều thread cùng muốn vào critical section. \texttt{OUTSIDE\_WORK}
    $= 0$ là tranh chấp tối đa.
  \item \emph{Vòng lặp đo.} Mỗi vòng, thread lấy lock, gọi \texttt{is\_prime(N)} trong critical
    section, nhả lock, rồi gọi \texttt{is\_prime(N)} thêm OW lần ở ngoài (OW viết tắt
    \texttt{OUTSIDE\_WORK}). $p$ là số thread. Mục Phương pháp đo định nghĩa đầy đủ các ký hiệu, độ
    chậm và độ công bằng.
  \item \emph{TAS, TTAS, MCS.} TAS lặp \texttt{exchange} lên cờ lock. TTAS đọc cờ tới khi rảnh rồi mới
    \texttt{exchange}~\cite{anderson1990spinlock}. MCS xếp hàng, mỗi thread chờ trên cờ
    riêng~\cite{mellorcrummey1991mcs}.
  \item \emph{Episode.} Một lần cả $p$ thread cùng vượt barrier~\cite{hensgen1988barrier}.
  \item \emph{FAA, CAS.} Fetch-and-add cộng và trả giá trị cũ trong một bước. Compare-and-swap chỉ ghi
    nếu giá trị chưa đổi, nếu không thì thất bại và phải thử lại.
\end{itemize}

Mục Kết quả giải thích số đo chủ yếu bằng số lần một cache line đi giữa các core trên đường tới
hạn~\cite{herlihy2020art}, cùng với scheduler và hai loại core của máy. Các hình đánh dấu hai mốc trên
trục số thread, là số P-core và số core của máy (trên \PCname{apple-m5} là \PCval{apple-m5}{pcores} và
\PCval{apple-m5}{ncpu}). P-core là core hiệu năng cao. E-core là core tiết kiệm điện, chậm hơn, và hai
loại nằm ở hai cụm có cache L2 riêng.

\begin{description}[leftmargin=0pt, itemsep=0.2em, font=\normalfont]
  \item[RQ1 (đúng đắn).] Các cài đặt có đúng không, và bộ kiểm tra có đủ nhạy để bắt lỗi cài đặt?
  \item[RQ2 (lock: mở rộng).] TAS, TTAS, MCS mở rộng thế nào theo số thread, \texttt{OUTSIDE\_WORK} và
    $N$, so với lock thư viện?
  \item[RQ3 (lock: backoff và công bằng).] Backoff giúp hay hại, và ở chế độ nào? Lock nhanh nhất có
    công bằng không?
  \item[RQ4 (barrier).] Chi phí một episode tăng thế nào theo số thread, theo cách chờ, so với barrier
    thư viện?
  \item[RQ5 (counter).] Thứ tự hiệu năng giữa FAA, CAS (có/không backoff) và counter dùng lock, và vì
    sao?
  \item[RQ6 (giữa các máy).] Các kết luận trên có giữ nguyên trên máy khác không?
\end{description}

\begin{table}[H]\centering\footnotesize
\begin{tabularx}{\linewidth}{@{}l>{\raggedright\arraybackslash}p{4.4cm}>{\raggedright\arraybackslash}X@{}}
\toprule
RQ & Thí nghiệm & Hình, bảng \\
\midrule
RQ1 & stress, TSan, mutation & Bảng~\ref{tab:corr:mutation}, \ref{tab:corr:summary} (phụ lục) \\
RQ2 & lock $\times$ thread $\times$ OW $\times$ $N$ &
  Hình~\ref{fig:apple-m5:lock_slowdown}, \ref{fig:apple-m5:lock_cs_length} (phụ lục);
  Bảng~\ref{tab:apple-m5:locktime}, \ref{tab:handoff} \\
RQ3 & sáu biến thể backoff, thời gian ngủ thật & Hình~\ref{fig:apple-m5:lock_backoff},
  \ref{fig:apple-m5:lock_fairness} \\
RQ4 & episode/giây, bảy barrier & Hình~\ref{fig:apple-m5:barrier_rate};
  Bảng~\ref{tab:apple-m5:barrier} \\
RQ5 & thời gian, số CAS hỏng & Hình~\ref{fig:apple-m5:counter_time},
  \ref{fig:apple-m5:counter_cas_failures}; Bảng~\ref{tab:apple-m5:countertime} \\
RQ6 & cùng quy trình trên \PCname{C4M7Y670NT} và \PCname{LAP1184} & Hình~\ref{fig:cross:lock},
  \ref{fig:cross:barrier}, \ref{fig:cross:counter}; Bảng~\ref{tab:cross:machines} và các bảng ở
  phụ lục \\
\bottomrule
\end{tabularx}
\caption{Câu hỏi nghiên cứu, thí nghiệm và hình, bảng trả lời.}\label{tab:intro:rqmap}
\end{table}

Đóng góp của báo cáo:
\begin{itemize}[nosep]
  \item cài đặt ba nhóm cấu trúc kèm lập luận memory order cho từng thao tác atomic;
  \item bộ kiểm tra (stress, watchdog, ThreadSanitizer, schedule fuzzing) được đánh giá bằng mutation
    testing;
  \item quy trình đo một lệnh cho đợt đo chính, sinh lại mọi hình và bảng của đợt đo đó, và thêm máy chỉ
    cần thêm tên máy vào danh sách của báo cáo;
  \item so sánh số đo với cận dưới lý thuyết.
\end{itemize}

Mục Nền tảng thí nghiệm mô tả máy đo, và mục Cài đặt trình bày thuật toán. Mục Kiểm tra tính đúng đắn
trả lời RQ1. Mục Phương pháp đo mô tả cách đo, và mục Kết quả và phân tích trả lời RQ2--RQ6. Báo cáo kết
thúc bằng mục Các yếu tố đe dọa độ tin cậy, mục Kết luận và phụ lục.

% ====================================================================================================
\section{Nền tảng thí nghiệm}
% TODO(report): 0,5 trang chữ + bảng.
%  - Đề yêu cầu số nền tảng shared-memory >= số thành viên: nêu rõ nhóm có mấy người, mấy máy.
%  - Với từng máy: vì sao nó đáng đo (kiến trúc x86 hay ARM, memory model TSO hay yếu hơn, số
%    core, có P-core/E-core hay SMT không, hệ điều hành có cho ghim thread không).
%  - Cache line: hệ thống báo bao nhiêu, mã nguồn đệm bao nhiêu, nhóm xác định bằng cách nào.
%  - OpenMP: máy nào có baseline OpenMP (compiler nhận -fopenmp), máy nào không.
%  - Điều kiện lúc đo: cắm sạc hay pin, chế độ tiết kiệm điện, các chương trình khác đã tắt chưa.
% Mục 2: Nền tảng thí nghiệm. Nguồn: results/apple-m5/machine.txt, results/apple-m5/campaign.log,
% include/platform.hpp, Makefile.

Đề yêu cầu số nền tảng shared-memory không ít hơn số thành viên. Báo cáo này có số đo của
\PCval{cross}{nhosts} máy, \PCname{apple-m5}. Mỗi thành viên thêm máy của mình bằng cách chạy
\texttt{make bench} trên máy rảnh, gửi lại thư mục \texttt{results/<tên máy>/} và thêm tên máy vào
danh sách máy của báo cáo. Bảng cấu hình (Bảng~\ref{tab:cross:machines}) và các hình của máy mới được
sinh lại tự động. Mọi nhận định về hiệu năng dưới đây chỉ nói về \PCname{apple-m5}.
\TODO{nhóm có mấy thành viên, và máy của các thành viên còn lại (sau khi họ chạy make bench)}

Máy \PCname{apple-m5} là một Apple M5 với 10 core không đồng nhất, gồm 4 P-core (performance) và 6
E-core (efficiency). Mỗi cụm có cache L2 riêng. Máy không có SMT. Ba đặc điểm của máy ảnh hưởng tới kết
quả. Thứ nhất, khi số thread vượt 4, một phần thread chạy trên E-core chậm hơn, và cache line dùng chung
phải đi qua lại giữa hai cụm. Thứ hai, macOS không cho gắn thread vào một core cố định, nên loại core
của mỗi thread do scheduler quyết định và có thể đổi giữa các lần chạy. Nhóm giảm biến thiên này bằng
cách lặp lại và lấy trung vị. Thứ ba, AArch64 có memory model yếu hơn x86-TSO. Các memory order khác
nhau sinh ra các lệnh máy khác nhau (store \texttt{relaxed} là \texttt{str}, store \texttt{release} là
\texttt{stlr}). Vì vậy lỗi chọn sai memory order có thể quan sát được trên máy này, trong khi trên x86 nó
thường không bao giờ lộ ra. CPU có các lệnh atomic LSE của ARMv8.1, nên \texttt{fetch\_add} và
\texttt{compare\_exchange} là một lệnh máy (\texttt{ldadd}, \texttt{cas}) chứ không phải vòng LL/SC.

Hệ điều hành báo cache line, đơn vị dữ liệu mà phần cứng giữ nhất quán giữa các core, dài 128 byte.
Mã nguồn đệm mọi biến dùng chung theo hằng \texttt{pc::kCacheLine}, bằng 128 trên Apple arm64 và 64
trên nền tảng khác. Nhóm kiểm tra rằng kích thước và căn chỉnh của mọi kiểu lock và của nút hàng đợi
MCS đều bằng 128 byte trên máy này.

Phần mềm gồm macOS 27.0.1 và Apple clang 21.0.0 với cờ \texttt{-std=c++20 -O2}. Apple clang không nhận
\texttt{-fopenmp}, và macOS không có \texttt{pthread\_spinlock\_t} hay \texttt{pthread\_barrier\_t}.
Vì vậy các baseline trên máy này là \texttt{std::mutex}, \texttt{pthread\_mutex\_t}, \texttt{os\_unfair\_lock},
\texttt{std::barrier} và barrier dựng từ mutex và condition variable.

Toàn bộ số đo đến từ một đợt đo duy nhất dài 1 giờ 26 phút, gồm 12\,152 lần chạy. Mỗi cấu hình có
7 lần lặp, và thứ tự các cấu hình xáo trộn trong từng lượt. Máy cắm sạc, chế độ tiết kiệm điện tắt, và
không có tiến trình đo hay build nào khác chạy cùng lúc. Load average một phút lúc bắt đầu là 6,73,
trong khoảng thường thấy của máy khi chỉ có ứng dụng nền. Không lần chạy nào sai kết quả kiểm tra.
Có 9 lần vượt thời hạn, đều là cấu hình có MCS khi số thread vượt số core, và chúng được đo lại với số vòng
lặp nhỏ hơn 100 lần (mục Phương pháp đo).

\PCtable{cross}{machines}

% ====================================================================================================
\section{Cài đặt}\label{sec:impl}
% TODO(report): với MỖI thuật toán: (1) đặc tả (bất biến, điều kiện trước/sau); (2) memory order
% của từng thao tác atomic và lý do, đặc biệt mọi chỗ dùng relaxed; (3) bố trí bộ nhớ (biến nào
% nằm riêng một cache line, vì sao); (4) giới hạn đã biết. Chỉ trích đoạn mã then chốt bằng
% môi trường lstlisting, không dán cả file. Trích nguồn cho thuật toán lấy từ sách và bài báo.
% Nguồn để viết: docs/part1.md, docs/part2.md, docs/part3.md và chú thích trong include/.

\subsection{Mutex lock}
% TODO(report):
%  - Giao diện acquire()/release() không tham số như khung của đề.
%  - TAS: mỗi lần thử là một phép ghi (exchange), kéo cache line về trạng thái độc quyền.
%  - TTAS: chờ bằng load, chỉ exchange khi thấy lock rảnh; vẫn có "bầy đàn" lúc nhả.
%  - Backoff: bốn chính sách của include/backoff.hpp (không backoff, ngủ đúng mẫu của đề, quay,
%    quay rồi yield); sleep_for là system call nên thời gian ngủ thật do hệ điều hành quyết định.
%  - MCS: hàng đợi, nút hàng đợi là biến thread_local; công bằng FIFO; giới hạn: một thread
%    không giữ được hai lock_mcs cùng lúc; vấn đề khi thread kế tiếp trong hàng bị preempt.
%  - Lock thư viện dùng để so sánh trên từng nền tảng (std::mutex, pthread_mutex_t,
%    os_unfair_lock, pthread_spinlock_t, omp_lock_t) và cách bọc chúng.
%  - Lệnh chờ trong vòng spin (pause / isb / yield) và lý do chọn.
% Mục 4.1. Nguồn: include/lock.h, backoff.hpp, platform.hpp. Lệnh AArch64: c++ -O2 -S (Apple clang 21).
Mọi lock trong \texttt{include/lock.h} có đúng \texttt{acquire()} và \texttt{release()} không tham số như
khung của đề. Với khai báo \texttt{using pc::lock\_mcs;}, dòng \texttt{lock\_mcs lock;} của khung đề
biên dịch nguyên văn. Lấy lock dùng acquire, nhả lock dùng release. Cặp này làm mọi phép ghi trong critical section trước nhìn thấy
được ở chủ lock sau. Trạng thái lock được khai báo \texttt{alignas(kCacheLine)}, nên mỗi lock chiếm riêng
một cache line (128 byte trên hai máy Apple, 64 byte trên \PCname{LAP1184}). Nhờ vậy không có false
sharing, tức hai biến độc lập nằm chung một line và làm mất hiệu lực line của nhau.
% nguồn: include/platform.hpp, kCacheLine = 128 khi __APPLE__ && __aarch64__, 64 ở nơi khác.

TAS lặp \texttt{exchange(true, acquire)} (lệnh \texttt{swpab}) tới khi đọc được \texttt{false}. TAS nhả lock bằng
\texttt{store(false, release)} (\texttt{stlrb}). TTAS chờ bằng \texttt{load(relaxed)} (\texttt{ldrb}) rồi
mới \texttt{exchange}. Phép đọc chỉ là gợi ý, \texttt{exchange} thắng mới là phép acquire.

Mọi vòng spin gọi \texttt{cpu\_relax()} giữa hai lần thử, tức là \texttt{pause} trên x86 và \texttt{isb}
trên AArch64. Lệnh \texttt{isb} được chọn vì \texttt{yield} của Apple gần như không làm gì (0,3~ns so với
12,1~ns của \texttt{isb}, cùng phép đo). \texttt{isb} dừng pipeline trong khoảng thời gian đó, nên
thread chờ đọc lock thưa hơn.
% nguồn: tables.tex, bảng sleeplatency: yield_insn trung vị 0,300 ns; isb trung vị 12,1 ns.

MCS~\cite{mellorcrummey1991mcs} giữ \texttt{tail\_} (\texttt{nullptr} là rảnh). Mỗi thread chờ trên cờ
\texttt{locked} trong nút của mình, nên mỗi lần chuyển giao chỉ ghi một line của một thread, theo thứ tự
FIFO. Trong các đoạn mã dưới đây, \texttt{relaxed}, \texttt{acquire}, \texttt{release},
\texttt{acq\_rel} viết tắt cho \texttt{std::memory\_order\_*}.
\begin{lstlisting}
me->next.store(nullptr, relaxed); me->locked.store(true, relaxed);  // not yet reachable
pred = tail_.exchange(me, acq_rel);                 // release: publish node; acquire: take lock
if (pred == nullptr) return;
pred->next.store(me, release);
while (me->locked.load(acquire)) cpu_relax();       // spin on own cache line
// release(): succ = next.load(acquire); if null, CAS tail_ me -> null (release);
//            on failure wait until next != null (acquire)
succ->locked.store(false, release);                 // hand over
\end{lstlisting}
Hai phép ghi relaxed an toàn vì nút chỉ trở nên với tới được qua hai phép release theo sau. Nếu CAS
\texttt{tail\_} thất bại, một thread đã đổi \texttt{tail\_} nhưng chưa nối vào \texttt{next}. Chủ lock
phải chờ nó, vì nếu bỏ đi thì thread đó chờ mãi. Vòng chờ này ngắn, vì thread kia đã \texttt{exchange}
xong và lệnh kế tiếp của nó là ghi \texttt{next}. Một thread chỉ vào critical section khi nó thấy
\texttt{tail\_} rỗng hoặc khi thread trước nó đặt \texttt{locked} của nó về \texttt{false}, nên tại mỗi
lúc chỉ có một chủ lock. Hàng đợi theo thứ tự FIFO, nên không thread nào bị
đói. Trên AArch64 các thao tác thành \texttt{swpal}
(\texttt{exchange} lên \texttt{tail\_}), \texttt{stlr} (nối \texttt{next}), \texttt{ldaprb} (chờ
\texttt{locked}), \texttt{ldapr} (\texttt{release()} đọc \texttt{next}), \texttt{casl} (CAS
\texttt{tail\_}) và \texttt{stlrb} (trao lock). Nút \PCname{mcs_node} cũng chiếm riêng một line.
% nguồn: c++ -O2 -S của mcs_acq/mcs_rel (Apple clang 21): swpal, stlr, ldaprb; ldapr, casl, stlrb.

Giới hạn thứ nhất của MCS là một thread không giữ được hai lock MCS cùng lúc, vì nút là biến
\texttt{thread\_local} dùng chung cho mọi \PCname{lock_mcs} do giao diện không có tham số. Giới hạn thứ
hai là lock được giao cho đúng một thread, nên nếu thread đó bị preempt thì mọi thread phía sau phải chờ
scheduler.

\PCname{lock_tas_bo<B>} và \PCname{lock_ttas_bo<B>} gọi chính sách \texttt{B} sau mỗi lần thử thất bại.
TTAS chỉ lùi lại khi thua \texttt{exchange} (BackoffLock~\cite{herlihy2020art}). Memory order giống TAS và
TTAS, nên lập luận về tính đúng ở trên giữ nguyên. \texttt{B} là một trong
ba chính sách. \PCname{sleep} theo mẫu của đề, gọi \texttt{sleep\_for} với \texttt{delay} nhân đôi từ
100~ns tới \texttt{MAX\_DELAY} $= 100\,000$~ns. \PCname{spin} quay ngẫu nhiên, với cửa sổ nhân đôi từ 4
tới 1024 lần \texttt{cpu\_relax}. \PCname{yield} quay 64 lần rồi gọi \texttt{yield}. \texttt{sleep\_for}
là system call. Trên \PCname{apple-m5}, yêu cầu 100~ns ngủ trung vị 2\,670~ns (nhóm đo bằng
\texttt{bench/sleep\_latency.cpp}).
% nguồn: tables.tex, bảng sleeplatency, sleep_for 100 ns: trung vị 2670 ns.

Lock thư viện được bọc inline sau cùng giao diện. Đó là \texttt{std::mutex}, \texttt{pthread\_mutex\_t},
\texttt{os\_unfair\_lock} (macOS), \texttt{pthread\_spinlock\_t} (Linux), \texttt{omp\_lock\_t} (khi có
OpenMP). Bảng~\ref{tab:impl:names} (phụ lục) liệt kê tên mọi cài đặt và máy có chúng.

\subsection{Sense-reversing barrier}
% TODO(report):
%  - Vì sao barrier một cờ không dùng lại được, và ý tưởng đảo chiều (sense reversal).
%  - Thứ tự hai thao tác của thread tới cuối (đặt lại bộ đếm rồi mới đảo sense) và vì sao.
%  - Các cách chờ: spin, backoff, chờ chặn; barrier thư viện dùng để so sánh.
%  - Memory order và bố trí bộ nhớ (local sense của mỗi thread nằm riêng một cache line).
% Mục 4.2. Nguồn: include/barrier.hpp; lệnh AArch64: c++ -O2 -S (Apple clang 21).
Barrier cho $n$ thread bảo đảm không thread nào ra khỏi lần gọi \texttt{wait()} thứ $k$ trước khi cả $n$
thread đã gọi lần thứ $k$. Barrier cũng dùng lại được qua mọi episode. Barrier một cờ ``mở/đóng'' không dùng lại
được. Đóng cờ sớm thì thread chậm chờ mãi, đóng muộn thì thread nhanh qua luôn episode sau.

Sense-reversing barrier~\cite{hensgen1988barrier} không đặt lại cờ mà đảo nó mỗi episode. Mỗi thread chờ
cờ bằng \emph{local sense} của mình, nên hai episode kề nhau chờ hai giá trị khác nhau. Giữa hai episode,
mọi local sense bằng \texttt{!sense\_}. Cờ không thể đảo lại trước khi thread đang chờ rời đi, vì lần
đảo sau cần đủ $n$ lần tới, trong đó có lần tới của chính thread này.
\begin{lstlisting}
const bool my_sense = local_[tid].sense;
if (count_.fetch_add(1, acq_rel) == n_ - 1) {      // last arriver
    count_.store(0, relaxed);                      // 1. reset for the next episode
    sense_.store(my_sense, release);               // 2. THEN release the others
} else {
    while (sense_.load(acquire) != my_sense) backoff.pause();
}
local_[tid].sense = !my_sense;
\end{lstlisting}
Thread tới cuối phải đặt lại bộ đếm trước khi đảo cờ. Nếu đảo trước, thread vừa được thả có thể tới
episode sau, rồi bị phép đặt lại muộn xoá mất lần tới, và mọi thread chờ mãi. Bộ kiểm tra bắt được biến thể đảo thứ
tự này (Bảng~\ref{tab:corr:mutation}).

\texttt{fetch\_add} cần acq\_rel (\texttt{ldaddal}). Nửa release công bố phép ghi trước barrier. Các
\texttt{fetch\_add} lên \texttt{count\_} nối nhau thành một release sequence, nên nửa acquire của thread
tới cuối đồng bộ với mọi thread tới trước. Phép ghi release lên \texttt{sense\_} (\texttt{stlrb}) chuyển
tiếp các phép ghi đó cho mọi thread đang chờ, vì chúng đọc cờ bằng acquire (\texttt{ldaprb}). Với chỉ
release (\texttt{ldaddl}), bài stress bắt được một thread qua barrier rồi đọc ô cũ của thread khác
(Bảng~\ref{tab:corr:mutation}). Phép đặt lại relaxed đủ vì nó đứng trước phép ghi release lên
\texttt{sense\_}.

\texttt{count\_}, \texttt{sense\_} và mỗi local sense nằm trên các cache line riêng, để lần tới không làm
mất hiệu lực line đang được chờ. Local sense nằm trong barrier, đánh chỉ số bằng \texttt{tid}, vì một biến
\texttt{thread\_local} sẽ bị hai barrier dùng xen kẽ làm hỏng.

Các cách chờ gồm quay (\PCname{sr}), ba chính sách backoff như lock, và chờ chặn (\PCname{sr_block}) bằng
\texttt{std::atomic::wait} và \texttt{notify\_all}. \PCname{sr_block} giữ nguyên giao thức và memory order
ở trên. \texttt{wait(old)} kiểm lại giá trị một cách atomic trước khi ngủ, nên không mất lần đánh thức nào.
Theo mã nguồn libc++ và libstdc++, \texttt{wait} quay một lúc rồi mới chặn thật, nên \PCname{sr_block}
thực chất là quay ngắn rồi ngủ.
% nguồn: include/barrier.hpp, chú thích của SenseBarrierBlocking ("the standard libraries poll for a short time").
Các barrier thư viện gồm \texttt{std::barrier},
barrier dựng từ mutex và condition variable (\PCname{condvar}), \texttt{pthread\_barrier\_t} (Linux) và
\texttt{omp barrier}. Các giới hạn là số thread cố định và barrier tập trung tốn $O(n)$ trên đường tới hạn. Biến thể quay
vẫn đúng khi số thread vượt số core, nhưng mỗi episode khi đó có thể tốn cả một lượt lập lịch.

\subsection{Concurrent counter}
% TODO(report):
%  - Ba loại đồng bộ đề yêu cầu: mutex lock (counter_lock<L> với từng lock của phần 1 và lock
%    thư viện), vòng lặp compare-and-swap (có và không có backoff), fetch-and-add.
%  - increment() trả về giá trị trước khi tăng; memory order (bản mặc định và bản seq_cst).
%  - Trên ARM compiler sinh lệnh atomic đơn (LSE) hay vòng LL/SC: xem assembly và nêu kết quả.
% Mục 4.3. Nguồn: include/counter.hpp. Lệnh máy: c++ -O2 -S (Apple clang 21), mặc định (LSE),
% -Xclang -target-feature -Xclang -lse, và -arch x86_64.
Counter bắt đầu từ 0. \texttt{increment()} cộng 1 và trả về giá trị \emph{trước} khi tăng. $M$ lần tăng
phải trả về đúng $0, \dots, M - 1$, mỗi giá trị một lần, và theo thứ tự thời gian thực
(linearizable~\cite{herlihy1990linearizability}). Ba cài đặt là \PCname{counter_lock<L>} (biến
\texttt{long} trong critical section của mỗi lock ở phần 1), \PCname{counter_cas} (lock-free, có biến thể
backoff) và \PCname{counter_faa} (wait-free). Trên \PCname{apple-m5} có \PCval{apple-m5}{ncounters} counter.
\begin{lstlisting}
long old = value_.load(relaxed);                   // a guess; the CAS checks it
while (!value_.compare_exchange_weak(old, old + 1, Order, relaxed)) {
    ++failed_;                                     // thread_local, failure path only
    backoff.pause();
}
return old;
// fetch-and-add: return value_.fetch_add(1, Order);
\end{lstlisting}

Thứ tự mặc định là \texttt{relaxed}, vì tính atomic của phép đọc-sửa-ghi và thứ tự toàn phần của mọi phép
ghi lên một biến atomic đúng với mọi memory order. Counter cũng không bảo vệ dữ liệu nào khác. Biến thể
\PCname{faa_seqcst}, \PCname{cas_seqcst} đo cái giá của \texttt{seq\_cst}. Counter dùng lock không có
atomic riêng. Nó đúng nhờ cặp release--acquire của lock.

Apple clang bật LSE (lệnh atomic đơn của ARMv8.1) mặc định. \texttt{fetch\_add} là một lệnh
\texttt{ldadd} (\texttt{ldaddal} với \texttt{seq\_cst}), không vòng lặp. CAS là một lệnh \texttt{cas}
(\texttt{casal}), chỉ hỏng khi giá trị thật sự đã đổi, nên mỗi lần \texttt{cas\_fail} là một xung đột thật.
Khi tắt LSE, cả hai thành \texttt{ldxr}/\texttt{stxr} (LL/SC), và CAS weak có thể hỏng giả. Trên x86-64
hai thứ tự sinh cùng một lệnh (\texttt{lock xadd}, \texttt{lock cmpxchg}).

Giá trị atomic nằm riêng một cache line. Trong \PCname{counter_lock}, lock và giá trị nằm trên hai line
để thread chờ không làm mất hiệu lực line của giá trị. Vì vậy mỗi lần tăng ghi lên hai line. Số lần CAS hỏng
được đếm trong biến \texttt{thread\_local} để không thêm một line bị tranh chấp. Về giới hạn, counter CAS
chỉ lock-free nên một thread có thể bị đói. \PCname{counter_lock} mang theo giới hạn của lock bên trong.

% ====================================================================================================
\section{Kiểm tra tính đúng đắn}
% TODO(report): chiến lược chung trước, rồi từng phần. Với mỗi bài kiểm tra: tính chất được kiểm,
% cách kiểm, kết quả trên TỪNG máy (make test, make tsan). Nêu rõ kiểm thử bằng stress chỉ tìm
% được lỗi chứ không chứng minh được là không có lỗi.

\subsection{Chiến lược}
% TODO(report): các tính chất cần kiểm (an toàn, sống còn, công bằng); stress test với số thread
% vượt số core; watchdog chống treo; ThreadSanitizer; mutation testing để đánh giá bộ kiểm tra.
% Mục 4.1: Chiến lược kiểm tra. Nguồn: tests/*.cpp, mk/part{1,2,3}.mk,
% results/{apple-m5,C4M7Y670NT,LAP1184}/{test,tsan,mutation}.log.

Nhóm kiểm hai nhóm tính chất. Tính an toàn (safety) nghĩa là không có hai thread
cùng ở trong critical section, không thread nào ra khỏi barrier sớm, không giá trị counter nào bị trả
về hai lần. Tính sống còn (liveness) nghĩa là chương trình không treo. Công bằng chỉ được đo.

Cùng một bài kiểm tra dạng template chạy cho mọi cài đặt, kể cả lock và barrier của thư viện.

Nhóm dùng bốn kỹ thuật:
\begin{itemize}[nosep]
  \item stress test với ít thread (1--4), với số luồng phần cứng và với gấp đôi số đó (\PCname{apple-m5}:
    10 và 20; \PCname{C4M7Y670NT} và \PCname{LAP1184}: 12 và 24); ở mức gấp đôi, thread bị tước core giữa giao thức;
  \item watchdog, thread giám sát báo \texttt{FAIL} khi bài kiểm tra ngừng tiến triển;
  \item ThreadSanitizer (TSan), công cụ của compiler báo hai truy cập xung đột, ít nhất một không atomic,
    mà không có quan hệ happens-before (atomic riêng của bài kiểm tra dùng \texttt{relaxed} để không tự tạo
    thứ tự, trừ đồng hồ \texttt{seq\_cst} của lượt kiểm thời gian thực);
  \item mutation testing, chạy bộ kiểm tra trên các biến thể cố ý làm hỏng.
\end{itemize}
Barrier dùng thêm schedule fuzzing. Counter dùng thêm phép kiểm thứ tự thời gian thực.

Bảng~\ref{tab:corr:summary} (phụ lục) tóm tắt kết quả trên ba máy. Mọi cài đặt đạt mọi bài kiểm tra, TSan
không báo data race nào trên mã đúng, và bộ kiểm tra bắt cả 14 biến thể hỏng. Bản TSan chạy với kích thước nhỏ hơn
(ít vòng hơn 10--20 lần) vì TSan làm chương trình chậm đi nhiều. Đầu ra nằm trong
\texttt{results/<máy>/}.
% nguồn: tests/test_locks.cpp (small: chia 10, 100 ms), tests/test_counter.cpp (small: chia 16),
% tests/test_barrier.cpp (small: g_size 0,05),
% mk/*.mk (make tsan chạy với đối số small).


\subsection{Mutex lock}
% TODO(report): loại trừ lẫn nhau (biến đếm không atomic + bộ đếm số thread đang ở trong critical
% section); không deadlock; công bằng (đo, không phán xét); quá tải; kết quả TSan.
\label{sec:corr:lock}
% Mục 4.2: Kiểm tra mutex lock. Nguồn: tests/test_locks.cpp, results/{apple-m5,C4M7Y670NT,LAP1184}/test.log, tsan.log.

Trong critical section, bài kiểm tra tăng một biến đếm không atomic (đọc, cộng một, ghi) và đếm số
thread đang ở bên trong. Vi phạm hiện ra trực tiếp, hoặc qua một lần cập nhật bị mất. Lượt đầu chạy
số vòng cố định: 1 thread lấy lock 200\,000 lần, 2 thread mỗi thread 100\,000 lần, và số luồng phần cứng
mỗi thread 20\,000 lần. Lượt sau chạy trong 300~ms với 4 thread, số luồng phần cứng và gấp đôi số đó
(\PCname{apple-m5}: 4, 10, 20; hai máy còn lại: 4, 12, 24). Lượt gấp đôi thử lock khi quá tải. Watchdog
kết thúc một bài kiểm tra con sau 20~s và báo deadlock.
% nguồn: tests/test_locks.cpp:220-232 (thread_counts {1, 2, hw}, iter_counts {200000, 100000, 20000},
% timed_threads {4, hw, 2hw}, ms = 300, watchdog 20 s).

Từ số lần lấy lock của từng thread, bài kiểm tra in tỉ số số lần lấy lock lớn nhất / nhỏ nhất (1,00 là
chia đều). Khi mỗi thread có một core, \PCname{mcs} chia gần như đều (1,00 ở 4 thread, 1,01 ở 10 thread),
đúng tính chất FIFO của hàng đợi MCS~\cite{mellorcrummey1991mcs}. Ở 20 thread, MCS chỉ lấy lock 15\,221
lần trong 300~ms (so với 2\,705\,564 lần ở 4 thread), với tỉ số 311,40. Số đo này phù hợp với việc lock
được trao cho một thread đang không có core (cơ chế bàn ở RQ2).
% nguồn: results/apple-m5/test.log, dòng mcs (một lần chạy 300 ms mỗi ô)

Trên \PCname{C4M7Y670NT} (12 core), MCS vẫn chia đều ở 12 thread (tỉ số 1,00, 1\,274\,328 lần lấy lock),
nhưng ở 24 thread chỉ lấy lock 10\,657 lần, với tỉ số 222,25. Trên cả hai máy Apple, MCS chỉ sụp khi số
thread vượt số core.
% nguồn: results/C4M7Y670NT/test.log, dòng mcs (4t 4\,354\,714 acq 1,00; 12t 1\,274\,328 acq 1,00;
% 24t 10\,657 acq 222,25).

Trên \PCname{LAP1184}, MCS sụp sớm hơn. Tỉ số là 1,00 ở 4 thread, nhưng MCS đã sụp ở 12 thread, bằng số
CPU ảo, với 10\,186 lần lấy lock (so với 3\,275\,092 ở 4 thread). Ở 24 thread chỉ còn 1\,628 lần. Ở
12 thread, nhóm cho rằng một CPU ảo đã bị Windows lấy core trong lúc giữ lượt của hàng đợi.
% nguồn: results/LAP1184/test.log, dòng mcs.

Ở bản TSan trên \PCname{LAP1184}, các lượt 100~ms với 24 thread không lấy được lock lần nào (\PCname{mcs}
và \PCname{ttas_bo_spin} ở 12 thread cũng 0). Vì vậy trên máy này TSan chỉ thực sự kiểm các lượt số vòng
cố định và các lượt 4--12 thread.
% nguồn: results/LAP1184/tsan.log, cột 24t = 0 acq ở cả 12 lock; mcs, ttas_bo_spin 12t = 0 acq.

Bản TSan trên \PCname{LAP1184} bỏ qua \PCname{omp_lock} và in ra việc bỏ qua. Khi không bỏ qua, TSan
báo một data race trên biến đếm mà \PCname{omp_lock} bảo vệ. Nhóm cho rằng đây là cảnh báo giả, vì
\texttt{omp\_set\_lock} và \texttt{omp\_unset\_lock} nằm trong libgomp, thư viện không được
instrument. Vì vậy TSan không thấy thứ tự mà lock tạo ra. Cùng lock đó đạt stress test ở bản thường.
% nguồn: results/LAP1184/tsan.log (dòng SKIP omp_lock); lần chạy TSan không bỏ qua (không lưu trong
% results/): SUMMARY: ThreadSanitizer: data race tests/test_locks.cpp:105 in critical_section.

\subsection{Sense-reversing barrier}
% TODO(report): không thread nào qua sớm; dùng lại qua rất nhiều episode liên tiếp; nhiễu thời
% điểm; TSan (lưu ý cảnh báo giả của std::barrier trên macOS, nếu có).
% Mục 4.3: Kiểm tra sense-reversing barrier. Nguồn: tests/test_barrier.cpp,
% results/apple-m5/test.log, tsan.log.

Ở episode $e$ (lần thứ $e$ cả nhóm đi qua barrier), mỗi thread ghi $e$ vào ô của mình, gọi
\texttt{wait}, rồi đọc mọi ô. Một ô nhỏ hơn $e$ nghĩa là thread này qua sớm. Một ô lớn hơn $e+1$ nghĩa
là thread khác đã qua episode kế tiếp mà không chờ. Các ô là atomic \texttt{relaxed}, nên điều kiện chỉ
đúng nhờ thứ tự do barrier tạo ra. Phép kiểm handoff dùng xen kẽ hai barrier với các ô là biến thường.
Dưới TSan, nó kiểm memory order của barrier. Các pha chạy 1, 4, 10 và 20 thread, tới 300\,000 episode liên tiếp. Watchdog báo treo khi không thread nào tiến trong 10~s.

Bản schedule fuzzing chèn một độ trễ ngẫu nhiên vào mỗi điểm giữa hai bước của giao thức, để một
thread bị chậm đúng giữa hai bước.

Trên \PCname{apple-m5}, cả 7 barrier đạt với tổng cộng 2\,198\,144 episode được kiểm. Bản fuzzing đạt
cho 5 barrier tự viết. Pha 20 thread của \PCname{sr} và \PCname{sr_bo_spin} hết ngân sách 1,5~s sau 64
và 80 episode. Dưới TSan cả 7 barrier đạt. Nhóm bỏ phép
kiểm handoff riêng cho \PCname{std_barrier}, vì libc++ của macOS không được instrument và TSan báo
cảnh báo giả.
% nguồn: results/apple-m5/test.log (cộng các dòng "PASS ... barrier episodes checked"), tsan.log

\subsection{Concurrent counter}
% TODO(report): giá trị cuối; mỗi giá trị trả về xuất hiện đúng một lần và dãy của từng thread
% tăng ngặt (bằng chứng linearizability cho counter); bẫy trình tối ưu che lỗi (docs/part3.md).
% Mục 4.4: Kiểm tra concurrent counter. Nguồn: tests/test_counter.cpp, results/apple-m5/test.log,
% tsan.log. Bẫy trình tối ưu: thí nghiệm phụ lúc phát triển (docs/part3.md mục 5.3, không trong results/).

\texttt{increment()} trả về giá trị trước khi tăng. Nhóm lưu mọi giá trị trả về và kiểm rằng tập giá trị trả về đúng là $\{0,\dots,M-1\}$. Dãy của
mỗi thread phải tăng ngặt. Ở một lượt riêng, nếu lần gọi A trả về trước khi B bắt đầu thì giá trị của A
phải nhỏ hơn. Các phép kiểm
này là bằng chứng linearizability~\cite{herlihy1990linearizability}. Tập giá trị đúng là
$\{0,\dots,M-1\}$ khi và chỉ khi có một thứ tự tuần tự hợp lệ. Hai phép kiểm thứ tự xác nhận thứ tự
đó tôn trọng thời gian thực.

Ở \texttt{-O2}, compiler gộp một vòng tăng không atomic thành một lần đọc và một lần ghi, nên bài kiểm
tra chỉ so giá trị cuối gần như luôn đạt với một counter không đồng bộ. Bài kiểm tra của nhóm lưu mọi
giá trị trả về và cho địa chỉ counter đi qua \texttt{pc::opaque()} ở mỗi vòng, buộc mỗi lần tăng đọc và
ghi thật.

Trên \PCname{apple-m5}, cả 19 counter đạt, tổng cộng hơn 121 triệu lần tăng được kiểm. Riêng
\PCname{lock_mcs} hết ngân sách thời gian ở 20 thread và chỉ được kiểm qua 1\,765\,312 lần tăng. Dưới TSan cả 19 counter đạt, không có cảnh báo.
% nguồn: results/apple-m5/test.log (tổng 121 573 312 cộng bằng python), tsan.log

\subsection{Mutation testing}
% TODO(report): bảng các biến thể cố ý làm hỏng của từng phần (make mutation): bài kiểm tra nào
% bắt được, sau bao lâu; biến thể nào chỉ ThreadSanitizer bắt được và vì sao.
% Mục 4.5: Mutation testing, trả lời RQ1. Nguồn: results/{apple-m5,C4M7Y670NT,LAP1184}/mutation.log (toàn bộ).
% Số episode lúc bắt lỗi: results/apple-m5/mutation_logs/barrier_mut2.run{2,5}.log,
% barrier_mut2_fuzz.run{1..5}.log; barrier_mut4.run1.log, barrier_mut4_fuzz.run1.log;
% results/C4M7Y670NT/mutation_logs/barrier_mut4.run{1,2,4,5}.log.

\label{sec:corr:mutation}
\begin{rqanswer}\textbf{Trả lời RQ1.} Không bài kiểm tra nào tìm thấy lỗi trên ba máy, và bộ kiểm tra bắt
cả 14 biến thể hỏng. Tính đúng dựa trên lập luận ở mục Cài đặt. Kiểm thử củng cố lập luận đó. Trên
\PCname{apple-m5}, \PCname{C4M7Y670NT} và \PCname{LAP1184}, mọi cài đặt đạt mọi bài kiểm tra (12 lock,
7 barrier và 19 counter trên hai máy Apple, 13, 9 và 20 trên \PCname{LAP1184}), và TSan không báo data
race nào. Hai lỗi memory order khó bắt nhất. Stress test không bắt \texttt{TAS\_RELEASE\_RELAXED} lần nào
(0/3 trên cả ba máy). Bản thường bắt \texttt{fetch\_add\_release\_only} 1/5 lần trên \PCname{apple-m5},
4/5 trên \PCname{C4M7Y670NT} và 0/5 trên \PCname{LAP1184}. TSan bắt cả hai lỗi trên mọi máy.\end{rqanswer}
% nguồn: results/{apple-m5,C4M7Y670NT,LAP1184}/mutation.log: TAS_RELEASE_RELAXED MISSED 0/3, TSan CAUGHT (cả ba);
% fetch_add_release_only plain caught 1/5 (apple-m5), 4/5 (C4M7Y670NT), SURVIVED 0/5 (LAP1184), DATA RACE reported (cả ba);
% results/{máy}/test.log, tsan.log (số cài đặt, 0 cảnh báo).

\emph{Bằng chứng.} Bảng~\ref{tab:corr:mutation} tóm tắt kết quả trên \PCname{apple-m5}. Biến thể lock là
macro sửa một chỗ trong \texttt{lock.h}. Biến thể barrier và counter là bản sao trong \texttt{tests/},
mỗi bản cài một lỗi. Riêng \texttt{single\_flag} và \texttt{block\_reserve} là thuật toán khác hẳn. Khi
chạy biến thể hỏng, watchdog báo treo sau 5~s với lock và 2~s với barrier, và mỗi pha barrier có ngân
sách 0,5~s.
% nguồn: include/lock.h (PC_MUTANT_TAS_*, PC_MUTANT_MCS_*), tests/test_barrier.cpp (namespace mutant, PC_MUTANT 1..5),
% tests/test_counter.cpp (PC_MUTANT_plain_long..block_reserve); mk/part1.mk:44 (watchdog=5);
% mk/part2.mk:63 (stall_s=2 cap_s=0.5).

\begin{table}[tbp]\centering\scriptsize\setlength{\tabcolsep}{3pt}
% nguồn: results/apple-m5/mutation.log
\begin{tabularx}{\linewidth}{@{}>{\ttfamily}l>{\raggedright\arraybackslash}Xlll@{}}
\toprule
\normalfont\textbf{Biến thể} & \textbf{Lỗi cài vào} & \textbf{Stress (+ fuzzing)} & \textbf{TSan} & \textbf{Bị bắt bởi} \\
\midrule
TAS\_CHECK\_THEN\_SET & \texttt{load} rồi \texttt{store} thay \texttt{exchange} & 3/3 vi phạm loại trừ & race & cả hai \\
TAS\_RELEASE\_RELAXED & nhả lock bằng \texttt{store} relaxed & 0/3 & race & chỉ TSan \\
MCS\_NO\_WAIT\_SUCCESSOR & nhả lock không chờ thread kế tiếp & 3/3 treo & không & stress \\
MCS\_ALWAYS\_CLEAR\_TAIL & luôn đặt \texttt{tail} về null & 3/3 treo & race & cả hai \\
\midrule
single\_flag & một cờ, không đảo chiều & 5/5 (5/5) sai episode & không & stress \\
flip\_before\_reset & đảo \texttt{sense} trước khi đặt lại bộ đếm & 2/5 (5/5) treo & không & stress \\
nonatomic\_count & \texttt{load}+\texttt{store} thay \texttt{fetch\_add} & 5/5 (5/5) treo & không & stress \\
fetch\_add\_release\_only & \texttt{fetch\_add} release thay acq\_rel & 1/5 (0/5) sai episode & race & chủ yếu TSan \\
thread\_local\_sense & local sense chung cho mọi barrier & 5/5 (5/5) sai giá trị handoff & race & cả hai \\
\midrule
plain\_long & \texttt{long} thường & 3/3 final, dup, miss, order & race & cả hai \\
load\_store & \texttt{load} rồi \texttt{store} atomic & 3/3 final, dup, miss, order & không & kiểm giá trị \\
cas\_no\_retry & CAS một lần, không thử lại & 3/3 final, dup, miss & không & kiểm giá trị \\
lock\_read\_outside & đọc giá trị trước \texttt{acquire()} & 3/3 final, dup, miss, order & race & cả hai \\
block\_reserve & giữ trước 8 giá trị mỗi thread & 3/3 chỉ real-time & không & kiểm giá trị \\
\bottomrule
\end{tabularx}
\caption{Mutation testing trên \PCname{apple-m5}: lock (trên), barrier (giữa), counter (dưới). Cột stress
cho số lần chạy bắt được trên số lần chạy (trong ngoặc: bản schedule fuzzing) cùng dấu hiệu hoặc phép kiểm
báo lỗi; cột TSan cho biết TSan có báo data race không.}
\label{tab:corr:mutation}
\end{table}

Trên \PCname{C4M7Y670NT}, mọi biến thể bị bắt bởi cùng kỹ thuật như trên \PCname{apple-m5}. Chỉ số lần
stress test bắt được khác. Ở bản thường, \texttt{fetch\_add\_release\_only} bị bắt 4/5 lần (so với
1/5) và \texttt{flip\_before\_reset} 1/5 lần (so với 2/5). Bản fuzzing cho cùng kết quả như trên
\PCname{apple-m5} (0/5 và 5/5). Kết quả TSan của máy này chỉ dựa trên dòng tóm tắt trong
\texttt{mutation.log}, vì log chi tiết của các lần chạy TSan không được lưu.
% nguồn: results/C4M7Y670NT/mutation.log (Part 2: flip_before_reset caught 1/5 | caught 5/5;
% fetch_add_release_only caught 4/5 | SURVIVED 0/5 | DATA RACE reported); các dòng còn lại trùng apple-m5.
% results/C4M7Y670NT/mutation_logs/ không có *_tsan.log, counter_muttsan_*, locks_muttsan_*.

Trên \PCname{LAP1184}, mọi biến thể bị bắt bởi cùng kỹ thuật như trên \PCname{apple-m5}, với ba khác
biệt. Thứ nhất, \texttt{fetch\_add\_release\_only} không bị stress test bắt lần nào (0/5, fuzzing 0/5),
chỉ TSan bắt. Thứ hai, bản thường bắt \texttt{flip\_before\_reset} 5/5 lần, so với 2/5 trên
\PCname{apple-m5}. Thứ ba, với \texttt{thread\_local\_sense}, bản TSan dừng ngay khi phép kiểm handoff
báo sai giá trị, và TSan không in cảnh báo nào trước đó. Trên máy này biến thể chỉ bị phép kiểm giá trị
bắt (5/5 ở bản thường và bản fuzzing).
% nguồn: results/LAP1184/mutation.log (Part 2: flip_before_reset caught 5/5 hang 5; fetch_add_release_only
% SURVIVED 0/5 | SURVIVED 0/5 | DATA RACE reported; thread_local_sense: no report; checks: wrong
% episode/value; Part 1: TAS_RELEASE_RELAXED MISSED 0/3, TSan CAUGHT);
% results/LAP1184/mutation_logs/barrier_mut5_tsan.log (FAIL [handoff]: WRONG VALUE, không có dòng TSan).

\emph{Vì sao.} Hai biến thể chỉ TSan bắt ổn định đều là lỗi memory order. \texttt{TAS\_RELEASE\_RELAXED}
đổi lệnh nhả lock \texttt{stlrb} (store release) thành \texttt{strb}. TSan bắt nó vì store relaxed không
tạo cạnh happens-before giữa hai chủ lock liên tiếp. Ba lần chạy stress không có cập nhật nào bị mất.
Nhóm chưa đo vì sao lỗi không lộ trên ARM. Một khả năng là các phép ghi trong critical section thường
đã tới cache trước khi phép ghi nhả lock được thấy, nên cửa sổ lỗi rất hẹp.
\texttt{fetch\_add\_release\_only} đổi \texttt{ldaddal} thành \texttt{ldaddl}, nên thread tới cuối có
thể đọc ô của thread khác trước khi \texttt{fetch\_add} của nó có hiệu lực. TSan bắt nó qua các ô thường
của phép kiểm handoff.

Trên x86-64, mọi \texttt{fetch\_add} đều là lệnh \texttt{lock xadd}. Mọi store, kể cả \texttt{relaxed}
và \texttt{release}, đều là \texttt{mov}. Vì vậy trên \PCname{LAP1184} hai biến thể memory order chạy đúng
lệnh máy của bản đúng, và không lịch chạy nào làm lộ lỗi.
% nguồn: biên dịch include/lock.h, tests/test_barrier.cpp (bản đúng và hai biến thể) bằng Apple clang
% -arch x86_64 -O2 -S: assembly của biến thể trùng bản đúng. LAP1184 dùng g++ 11.5, chưa so assembly riêng.

Ngược lại, TSan im lặng với bảy biến thể chỉ dùng atomic hoặc gây treo. Không kỹ thuật nào một mình
bắt được mọi biến thể. Stress test hoặc phép kiểm giá trị bắt mọi lỗi thuật toán. Bốn lỗi treo (hai biến
thể MCS, \texttt{flip\_before\_reset}, \texttt{nonatomic\_count}) hiện ra qua watchdog.

Fuzzing làm \texttt{flip\_before\_reset} treo sau 3--7 episode (bản thường cần 39\,896 và 289\,052
episode ở hai lần bắt được trên \PCname{apple-m5}). Bản fuzzing không bắt được lỗi thiếu acquire (0/5).
Bản này chỉ chạy một phần mười số episode, và trên hai máy Apple bản thường chỉ bắt được lỗi sau hơn
13\,000 episode của một pha. Nhóm chưa đo vì sao bản thường trên \PCname{LAP1184} bắt
\texttt{flip\_before\_reset} thường hơn. Một khả năng là máy ảo tước core của thread nhiều hơn.
% nguồn: results/apple-m5/mutation_logs/barrier_mut2.run{2,5}.log (39 896, 289 052),
% barrier_mut2_fuzz.run{1..5}.log (5, 7, 7, 3, 5 episode); mutation.log (fetch_add_release_only fuzzing 0/5).
% Episode lúc bắt fetch_add_release_only ở bản thường: apple-m5 barrier_mut4.run1.log 35 881 (pha reuse);
% C4M7Y670NT barrier_mut4.run{1,2,4,5}.log 31 763, 13 373, 23 689, 26 190 (pha all-cores).
% Bản fuzzing: g_size = 0.1 (tests/test_barrier.cpp, main), pha reuse 30 000 và all-cores 10 000 episode
% (apple-m5/mutation_logs/barrier_mut4_fuzz.run1.log).

\emph{Độ vững.} Kết luận trên dựa trên một lượt chạy mỗi máy và chỉ trên các lịch chạy đã xảy ra.
Mục~\ref{sec:corr:limits} nêu các giới hạn này.

\subsection{Giới hạn của phương pháp}
% TODO(report): những gì bộ kiểm tra không chứng minh; lỗi memory order không lộ trên x86 (TSO);
% những nền tảng chưa chạy được TSan.
% Mục 4.6: Giới hạn của phương pháp kiểm tra (giới hạn của việc đo nằm ở mục Các yếu tố đe dọa).

Stress test và fuzzing chỉ xét các lịch chạy đã xảy ra. Chúng chỉ ra được là có lỗi, không chứng minh
được là không có lỗi. Tính đúng vì vậy dựa trên lập luận ở mục Cài đặt. Trên x86-64 (memory model TSO)
các lỗi memory order sinh cùng lệnh máy với bản đúng nên không bao giờ lộ ra lúc chạy. Với loại lỗi
này, nhóm chỉ có TSan. TSan chỉ tìm data race, chỉ thấy các đường đi đã chạy và cần mã được instrument.
Các baseline OpenMP và các lock, barrier chỉ có trên Linux chưa chạy dưới TSan. Bộ kiểm tra không kiểm
tiến triển (lock-free, wait-free) ngoài trường hợp treo, và không kiểm \texttt{get()} trong lúc các
thread đang tăng. Các biến thể đột biến do chính nhóm chọn, nên chúng chỉ chứng tỏ bộ kiểm tra nhạy
với những lỗi đó.

% ====================================================================================================
\section{Phương pháp đo}
% TODO(report): 1-1,5 trang. Mọi chi tiết dưới đây lấy từ scripts/run_bench.py (khối "THE GRID"
% và chú thích đầu file) và include/bench_common.hpp; đừng chép số từ trí nhớ.
%  - Vùng đo: từ lúc mở cổng xuất phát tới khi thread cuối cùng join xong; tạo thread nằm ngoài.
%  - Chống trình tối ưu xoá việc: is_prime không inline, N qua pc::opaque() ở mỗi lời gọi, tổng
%    kết quả được kiểm ở cuối (cột ok). Bằng chứng: bảng "đơn vị công việc" bên dưới (t gần như
%    không đổi theo OUTSIDE_WORK) và assembly.
%  - Lưới tham số: số thread (kể cả vượt số core), OUTSIDE_WORK, N; vì sao chọn các giá trị đó.
%  - Số vòng lặp: lock dùng I = W_N / max(p, 1 + OW) để cận dưới của mọi lần chạy dài như nhau;
%    barrier và counter hiệu chỉnh I cho từng cấu hình bằng một lần chạy thử. Vì I khác nhau giữa
%    các cấu hình, hình và bảng so sánh qua thời gian quy về cùng số vòng lặp.
%  - Cấu hình sụp đổ (lock hàng đợi khi số thread vượt số core): đo với I nhỏ hơn 100 lần.
%  - Thống kê: số lần lặp; báo cáo trung vị, min, max; vì sao không dùng trung bình.
%  - Thứ tự chạy: lặp là vòng ngoài, các cấu hình xáo trộn trong từng lượt lặp với seed cố định.
%  - Kiểm soát nhiễu: một tiến trình đo tại một thời điểm; ngưỡng load average lúc bắt đầu; load
%    ghi trong từng dòng số liệu; không ghim thread (macOS không cho).
%  - Thời hạn của một lần chạy và cách ghi nhận lần chạy vượt thời hạn (không bỏ, không coi là 0).
%  - Mô hình so sánh của phần lock: cận dưới max(p I t, (1 + OW) I t), t suy ra từ lần chạy 1 thread.
% Mục 5: Phương pháp đo. Nguồn: scripts/run_bench.py (chú thích đầu file, khối "THE GRID"),
% include/bench_common.hpp, include/workload.hpp, scripts/plot.py, results/apple-m5/{calibration.json,
% campaign.log, *.csv}, tables.tex (campaign, lockunit). Số tự tính: scripts/report_numbers/corr_stats.py.
% Giới hạn của việc đo chỉ nằm ở mục Các yếu tố đe dọa độ tin cậy. Bản dài: agent-corr/long/method.tex.

Mỗi lần chạy in đúng một dòng CSV gồm tham số, thời gian, cột \texttt{ok} (phép kiểm của lần
chạy) và load average. Đồng hồ chạy từ lúc mọi thread
đã được tạo và cổng xuất phát chung mở ra tới khi lần \texttt{join} cuối trả về. Việc tạo thread nằm
ngoài vùng đo.

Trong khung của đề, compiler được phép xoá \texttt{is\_prime(N)} vì kết quả bị bỏ đi. Nhóm giữ
\texttt{is\_prime} không inline, cho \texttt{N} đi qua \texttt{pc::opaque()} ở mỗi lời gọi, và kiểm tổng các kết quả ở cuối lần chạy (cột \texttt{ok}).
Bảng~\ref{tab:apple-m5:lockunit} xác nhận phần việc không bị xoá. Ở $N=1009$, thời gian một lời gọi
$t$ chỉ thay đổi từ 504 tới 527~ns khi \texttt{OUTSIDE\_WORK} tăng từ 0 lên 15, và $t$ tỉ lệ với $N$.
Cả 12\,152 lần chạy đều có \texttt{ok = 1}.
% nguồn: tables.tex, bảng lockunit; campaign.log: bad_check=0

Phần lock quét 12 lock qua ba lượt. Một lượt dùng $N=1009$ với \texttt{OUTSIDE\_WORK} (OW) bằng 0, 1, 3, 7, 15. Hai lượt còn lại dùng $N=11$
và $N=10007$ với OW bằng 0 và 3. $N$ phải là số nguyên tố để \texttt{is\_prime} chạy đủ $N-2$ vòng. Ba
giá trị cho critical section khoảng 5~ns, 0,5~\textmu s và 5~\textmu s. OW quyết định mức tranh chấp.
Theo mô hình, lock bão hoà khi $p > 1+\text{OW}$. Barrier quét 7 cài đặt. Counter quét 19 cài đặt với OW bằng 0, 1, 8
và $N=101$. Số thread trên \PCname{apple-m5} là 1, 2, 3, 4, 6, 8, 10, 12, 16, 20. Trong đó 4 là số P-core, 10 là
mọi core, 12--20 là quá tải.

Với lock, số vòng mỗi thread là $I = W_N/\max(p,\,1+\text{OW})$, để cận dưới của mọi cấu hình khoảng
0,19~s. Nhóm hiệu chỉnh $W_N$ một lần cho mỗi $N$ bằng \PCname{ttas}. Với barrier và counter, chi phí
một thao tác chênh nhau nhiều bậc, nên mỗi cấu hình có $I$ riêng suy từ một lần chạy thử, nhắm 0,25~s
và 0,15~s. $I$ cố định cho cả đợt đo. Hình và bảng quy thời gian về cùng số vòng lặp.

Thời hạn mỗi lần chạy là 30~s cho lock và 10~s cho barrier và counter. Nhóm đo lock MCS quá tải với số
vòng nhỏ hơn 100 lần ngay từ đầu (27 cấu hình). Lần chạy vượt thời hạn được ghi thành một dòng
\texttt{ok=0}, không bỏ đi, rồi cấu hình đó chạy lại với số vòng nhỏ hơn 100 lần. Trên
\PCname{apple-m5} có 9 lần như vậy, đều là MCS quá tải (6 ở lock, 3 ở counter).
% nguồn: lock.csv, counter.csv (dòng extra=timeout); calibration.json (khoá reduced|...)

Mỗi cấu hình chạy 7 lần. Hình báo cáo trung vị, thanh sai số đi từ min tới max. Nhóm không dùng trung
bình vì phân bố lệch phải. Ở 33\,\% cấu hình lock, lần chậm nhất dài hơn 1,5 lần trung vị. Lần vượt
thời hạn được tính là chậm nhất khi lấy trung vị. Trong mỗi lượt lặp, mọi cấu hình chạy theo thứ tự xáo
trộn với seed cố định. Mỗi thời điểm chỉ
có một tiến trình đo, và đợt đo chỉ bắt đầu khi load average không quá 7,0.
% nguồn: agent-corr/stats.py trên lock.csv (ok=1): tỉ lệ cấu hình có max/median > 1,5 = 0,334

Phần lock được so với cận dưới $T_{\min} = \max\bigl(p\,I\,t,\ (1+\text{OW})\,I\,t\bigr)$. Vế đầu là
tổng các critical section chạy nối tiếp, vế sau là phần việc của riêng một thread. $t$ là trung vị lần
chạy 1 thread của lock nhanh nhất chia cho $I(1+\text{OW})$. Thời gian đo chia cho $T_{\min}$ là hệ số
chậm. Bảng~\ref{tab:apple-m5:campaign} tóm tắt đợt đo. Bảng~\ref{tab:apple-m5:lockunit} cho $t$.


% ====================================================================================================
\section{Kết quả và phân tích}
% TODO(report): với MỖI hình viết đủ ba ý:
%   (1) Quan sát: hình cho thấy gì, có số ("ở 8 thread, TAS chậm hơn TTAS x lần");
%   (2) Giải thích bằng cơ chế: lưu lượng coherence, tranh chấp, công bằng, scheduler, P/E-core;
%   (3) Bằng chứng cho giải thích: số đo phụ nào ủng hộ (số lần CAS hỏng, độ công bằng, thí
%       nghiệm đối chứng chỉ đổi một yếu tố).
% Nói rõ kết luận nào được số liệu hỗ trợ và kết luận nào chỉ là giả thuyết. Mọi con số trích trong lời văn phải lấy từ bảng do
% plot.py sinh ra, không tự tính tay.

\subsection{Mutex lock}
% TODO(report):
%  - Thời gian và thông lượng theo số thread: lock nào bám sát cận dưới, tới đâu; đường cong gãy
%    ở số P-core, ở tổng số core hay ở điểm nào khác; chuyện gì xảy ra khi số thread vượt số core
%    (đây là kết quả, không phải lỗi đo) và vì sao lock thư viện hành xử khác (block thay vì spin).
%  - Ảnh hưởng của OUTSIDE_WORK: lock bão hoà khi p > 1 + OW.
%  - Ảnh hưởng của N (độ dài critical section): chi phí chuyển giao lock so với phần việc.
%  - Backoff: giúp hay hại, ở chế độ nào; giải thích bằng thời gian ngủ thật (bảng sleep_latency).
%  - Công bằng: lock "nhanh nhất" có công bằng không.
%  - Chi phí khi không tranh chấp (1 thread).
% Mục Kết quả, mutex lock (RQ2, RQ3). Nguồn: results/apple-m5/lock.csv, lock_yield.csv,
% sleep_latency.csv (trung vị 7 lần lặp); bảng locktime, locktime-ow*, locksingle, sleeplatency,
% lockunit trong report/figures/apple-m5/tables.tex. Số không có trong bảng in ra (độ chậm, độ công
% bằng, ns mỗi critical section, tỉ số backoff) tính bằng scripts/report_numbers/rlock_analyze.py và
% scripts/report_numbers/rlock_analyze2.py; hai script nạp scripts/plot.py và dùng đúng aggregate/unit_time/bound của nó, nên trùng
% với số vẽ trên hình. Số ghi "lock.csv: ..." tính thẳng từ CSV: ns mỗi critical section = trung vị
% seconds/ops*1e9, độ công bằng = trung vị thr_min_s/thr_max_s của cột extra. Nguồn từng số ghi trong comment ngay cạnh.

Hai đại lượng dùng trong mục này là độ chậm và độ công bằng, định nghĩa ở mục Phương pháp đo.

% ================================================================================================
\subsubsection{RQ2: Lock mở rộng thế nào theo số thread, tranh chấp và độ dài critical section?}

\begin{rqanswer}\textbf{Trả lời RQ2.} TTAS mở rộng tốt nhất trong ba lock tự viết. Ở 20 thread
($N = 1009$, $\mathrm{OW} = 0$), nó chậm 2,48 lần cận dưới. MCS gãy mạnh tại số core (13 lần ở 10
thread, 7 lần khi tải nền thấp) và sụp khi vượt số core (1\,700--6\,900 lần). Các ô đó đo với số vòng lặp giảm 100 lần và chỉ cho bậc độ
lớn. Lock thư viện không sụp. \PCname{std_mutex} và \PCname{pthread_mutex} chậm hơn spinlock khi
critical section dài cỡ 0,5~$\mu$s (3,9--4,0 lần cận dưới ở 20 thread, so với 2,48 của TTAS). Riêng
\PCname{os_unfair_lock} giữ 1,6--1,7 lần. Khi critical section chỉ vài nanosecond, cả ba lock thư viện
nhanh hơn spinlock không backoff từ 3 thread. Số liệu của \PCname{apple-m5}. RQ6 kiểm lại trên hai máy còn lại.\end{rqanswer}
% nguồn: scripts/report_numbers/rlock_analyze.py, độ chậm N=1009 OW=0: p=20 ttas 2,48, std_mutex 3,93, pthread_mutex 3,97,
% os_unfair_lock (p=2..20) 1,59..1,69; mcs p=10 13,0, p=12/16/20 1,70e3/5,16e3/6,88e3 (đo giảm vòng).
% N=11 OW=0 p>=3: std_mutex 4,08..5,81, pthread_mutex 3,68..5,14, os_unfair_lock 2,87..5,02; tas 5,18..80,5, ttas 7,28..28,2,
% mcs 18,7..101 (p<=10). Ở p=2 os_unfair_lock 5,57 > tas 4,39, nên câu chỉ đúng từ 3 thread.

\PCfigure{apple-m5}{lock_slowdown}

\emph{Bằng chứng.} Ở $\mathrm{OW} = 0$, ba spinlock tự viết chậm không quá 1,24 lần cận dưới tới 4
thread (Hình~\ref{fig:apple-m5:lock_slowdown}). Chúng gãy từ 4 lên 6 thread (TAS 1,18 lên 1,50, TTAS
1,08 lên 1,31, MCS 1,24 lên 2,98), tức tại số P-core chứ không tại số core. MCS gãy mạnh nhất. Từ 6
thread nó chậm nhất trong ba spinlock tự viết, và từ 8 thread chậm nhất trong mọi lock (4,31 so với 3,72
của \PCname{std_mutex}). Nó chậm 13 lần ở đúng 10 thread rồi sụp hẳn khi vượt số core
(Bảng~\ref{tab:apple-m5:locktime}). TAS và TTAS chỉ chậm dần, tới 3,48 và 2,48 lần ở 20 thread. Lock
thư viện đã chậm 2,5--2,9 lần từ 3--4 thread (\PCname{os_unfair_lock} 1,6 lần từ 2 thread), rồi tăng
chậm và gần như phẳng từ 8 thread (3,5--4,0 lần).
% nguồn: scripts/report_numbers/rlock_analyze.py, độ chậm N=1009 OW=0 p=4/6: tas 1,18/1,50, ttas 1,08/1,31, mcs 1,24/2,98;
% mcs p=8 4,31, std_mutex p=8 3,72; mcs p=10 13,0; p=20 tas 3,48, ttas 2,48.
% std_mutex p=3/4 2,79/2,93; pthread_mutex 2,45/2,86; p=8..20 std_mutex 3,72..3,93, pthread_mutex 3,54..3,97;
% os_unfair_lock (p=2..20) 1,59..1,69.

\PCtable{apple-m5}{locktime}

Ô $\mathrm{OW} = 15$ là một thí nghiệm đối chứng, vì với $p \le 15$ lock bận tối đa $p/16$ thời gian.
Vậy mà ở 6 thread cả mười hai lock, kể cả lock thư viện, cùng chậm 1,53--2,03 lần cận dưới.
% nguồn: scripts/report_numbers/rlock_analyze.py, độ chậm N=1009 OW=15 p=6: 1,53 (mcs) .. 2,03 (os_unfair_lock).

Từ $p = 1+\mathrm{OW}$, lock phải bận liên tục thì mới đạt cận dưới, nên mỗi lần chuyển giao chậm
đều cộng thẳng vào thời gian chạy. Lock thư viện cho thấy điểm này rõ nhất. Ở $\mathrm{OW} = 1$ và 2
thread, \PCname{std_mutex} và \PCname{pthread_mutex} đã chậm 2,9 lần cận dưới (\PCname{os_unfair_lock}
1,7 lần), còn ba spinlock tự viết chậm 1,11--1,27 lần. Ở $\mathrm{OW} = 3$, ba lock thư viện chậm
1,34--1,83 lần ở 3 thread và 2,64--2,86 lần ở 4 thread, còn spinlock 1,11--1,13 và 1,31--1,37 lần. Nhóm cho rằng
thread chờ của lock thư viện ngủ, nên mỗi lần chuyển giao trả thêm một lần đánh thức qua kernel (ước
lượng ở dưới), và lock thư viện bão hoà sớm hơn spinlock. Với $\mathrm{OW} = 7$ và 15, chỗ gãy từ 4 lên
6 thread tới trước điểm bão hoà (TAS 1,09 lên 1,92 và 1,05 lên 1,59), và mọi lock tăng cùng nhau.
% nguồn: scripts/report_numbers/rlock_analyze.py, độ chậm N=1009 OW=1 p=2: std_mutex 2,86, pthread_mutex 2,91,
% os_unfair_lock 1,71; tas/ttas/mcs 1,12/1,11/1,27. OW=3 p=3: std_mutex 1,83, pthread_mutex 1,68, os_unfair_lock 1,34,
% tas/ttas/mcs 1,13/1,11/1,12; p=4: std_mutex 2,64, pthread_mutex 2,86, os_unfair_lock 2,84, tas/ttas/mcs 1,31/1,37/1,37.
% tas OW=7 p=4/6 1,09/1,92; OW=15 p=4/6 1,05/1,59.

Khi vượt số core ở $\mathrm{OW} = 15$, riêng phần việc đã vượt sức máy. Tổng $p(1+\mathrm{OW})\,I$ lời
gọi \texttt{is\_prime} chia cho 10 core cần ít nhất $p(1+\mathrm{OW})\,I\,t/10$, tức 1,6 lần cận dưới ở
16 và 20 thread, trước khi tính chi phí lock. TAS, TTAS và lock thư viện cùng nằm trong 4,9--6,0 lần ở
16--20 thread. Nhóm cho rằng phần còn lại đến từ E-core chậm hơn và từ việc chia lượt lập lịch.
% nguồn: số học từ định nghĩa cận dưới max(p,1+OW) I t (mục Phương pháp đo): p(1+OW)/(10 max(p,1+OW)) = 16*16/160 = 1,6
% (p=16), 20*16/200 = 1,6 (p=20). scripts/report_numbers/rlock_analyze.py, độ chậm N=1009 OW=15 p=16/20: tas 5,28/6,02,
% ttas 4,93/5,82, std_mutex 4,92/4,95, pthread_mutex 4,88/5,07, os_unfair_lock 5,07/5,06.

Với $N = 10\,007$ (khoảng 5,2~$\mu$s) TAS và TTAS gần như trùng nhau
(Hình~\ref{fig:apple-m5:lock_cs_length}, phụ lục). Với $N = 11$ (4,6~ns), ở 20 thread TAS chậm 80,5 lần cận
dưới, TTAS 28,2 lần, còn lock thư viện chỉ khoảng 5 lần.
% nguồn: scripts/report_numbers/rlock_analyze.py, độ chậm OW=0: N=10007 p=20 tas 2,55, ttas 2,53; N=11 p=20 tas 80,5, ttas 28,2,
% std_mutex 5,71, pthread_mutex 4,97, os_unfair_lock 5,02.

\emph{Vì sao.} Nhóm cho rằng chỗ gãy ở 4 thread đến từ máy, không từ thuật toán. Ở ô đối chứng, mọi
lock cùng chậm như nhau. Độ công bằng ở 6 thread là 0,84--0,98, tức mọi thread cùng chậm đi. Điều này
phù hợp với việc scheduler luân chuyển thread giữa P-core và E-core chậm hơn, hoặc với tải nền (xem mục
Các yếu tố đe dọa độ tin cậy). Trên \PCname{C4M7Y670NT} (6 P-core, 12 core), điểm gãy dời sang 6 lên 8
thread và MCS sụp khi vượt 12 thread (mục So sánh giữa các máy).
% nguồn: scripts/report_numbers/rlock_analyze.py, fairness N=1009 OW=15 p=6: 0,84 (mcs) .. 0,98.
% nguồn: results/C4M7Y670NT/lock.csv (độ chậm tính như rlock_analyze.py: trung vị / (iters t p), t từ lần chạy 1 thread
% nhanh nhất), N=1009 OW=0: tas p=6/8 1,15/1,53; mcs p=6/8 1,23/2,58 (bảng locktime: 1,08/0,937, 1,15/0,937), p=12 4,99, p=16 2\,070.

MCS gãy mạnh ngay tại số core và sụp khi vượt số core, sớm hơn mức mà riêng lưu lượng cache coherence
giải thích được. Nguyên nhân nằm ở cách trao lock. MCS trao lock cho đúng một thread kế tiếp theo thứ tự FIFO (vào trước được phục
vụ trước). Nếu thread đó không có core, lock nằm yên cho tới khi scheduler cho nó chạy lại, và cả hàng
chờ mất một lượt lập lịch. TAS và TTAS không bị như vậy, vì thread nào đang chạy cũng lấy được lock
rảnh. Độ công bằng của MCS bằng 1 tới 10 thread (Hình~\ref{fig:apple-m5:lock_fairness}), tức lock luôn
đi đúng thứ tự hàng đợi.
% nguồn: scripts/report_numbers/rlock_analyze.py, fairness mcs N=1009 OW=0 p=2..10: 1,00.

Chương trình \texttt{handoff\_gap} kiểm cơ chế này. Nó chạy đúng vòng lặp của phần lock ($N = 1009$,
$\mathrm{OW} = 0$) cho MCS, TAS, TTAS và \PCname{std_mutex} ở 4, 8, 10 và 12 thread. Mỗi lần chạy có
khoảng 40\,000 critical section (3\,000 cho MCS ở 12 thread), và mỗi cấu hình chạy 5 lần theo thứ tự
xáo trộn. Thread nhả lock ghi lại thời điểm nhả, và thread lấy lock kế tiếp đo \emph{khoảng trống trao
lock}, tức thời gian từ lúc lock được nhả tới lúc nó có chủ mới. Ở 4 thread, trung vị của khoảng trống
dài nhất trong một lần chạy là 6--22~$\mu$s với cả bốn lock. Vì vậy một khoảng trống dài hơn
100~$\mu$s được coi là lock đã chờ một thread không có core. Ở 4 thread, chỉ một trong 20 lần chạy có
khoảng trống như vậy (MCS, hai lần trao lock).
% nguồn: experiments/run_extra.py (part_a: total = 3000 nếu mcs và p>=12, ngược lại 40000; iters = total // p;
% 5 rep, random.Random(SEED+rep).shuffle). handoff.csv: handovers 39\,999 (p=4,8,10), 39\,995 (p=12), 2\,999 (mcs p=12).
% max_gap_us p=4, trung vị 5 lần (extra_analyze.py): ttas 6, tas 10, std_mutex 14, mcs 22; lần lớn nhất: mcs 199,5
% (rep 3, long_handovers 2), tas 79,1, ttas 69,0, std_mutex 18,1.

Bảng~\ref{tab:handoff} cho thấy MCS ở 10 thread chỉ có 0,25\,\% lần trao lock dài như vậy, mỗi lần
trung bình 0,79~ms, nhưng chúng chiếm 47\,\% thời gian chạy. Ở 12 thread, 27\,\% lần trao lock là lần
dài và chúng chiếm 99\,\% thời gian. TAS, TTAS và \PCname{std_mutex} gần như không có lần trao lock dài
nào. Số lần chuyển ngữ cảnh không tự nguyện của tiến trình MCS tăng từ 60 ở 4 thread lên 1\,970 ở 10
thread và 23\,044 ở 12 thread. Phép đo này chạy khi load average khoảng 2,8, nên ngay ở 10 thread đã có
thread bị tước core. Đối chứng là TAS ở 10 thread. Tiến trình TAS cũng bị chuyển ngữ cảnh 1\,022 lần,
cùng bậc với MCS, nhưng trung vị số lần trao lock dài bằng 0, vì thread nào đang chạy cũng lấy được lock
rảnh. Phép đo này không kiểm điểm gãy từ 4 lên 6 thread, nên cơ chế của điểm gãy đó vẫn là giả thuyết.
% nguồn: scripts/report_numbers/extra_analyze.py (results/apple-m5/extra/handoff.csv, trung vị 5 lần chạy):
% mcs p=10 long 0,245 %, mean_long 792 us, time_in_long 47,3 %; p=12 26,8 % / 2\,233 us / 99,4 %;
% nivcsw mcs p=4/10/12: 60 / 1\,970 / 23\,044; load1 2,8--2,9.
% tas p=10: nivcsw 1022/2317/466/753/1024 (trung vị 1\,022), long_handovers 0/8/2/0/0 (trung vị 0).

\begin{table}[H]\centering\small
\begin{tabular}{@{}lcccc@{}}
\toprule
 & 4 thread & 8 thread & 10 thread & 12 thread \\
\midrule
\PCname{mcs}       & 0 / 0    & 0,057 / 8,3 & 0,25 / 47 & 27 / 99 \\
\PCname{tas}       & 0 / 0    & 0 / 0       & 0 / 0     & 0,003 / 0,3 \\
\PCname{ttas}      & 0 / 0    & 0 / 0       & 0 / 0     & 0 / 0 \\
\PCname{std_mutex} & 0 / 0    & 0 / 0       & 0 / 0     & 0 / 0 \\
\bottomrule
\end{tabular}
\caption{Lần trao lock dài hơn 100~$\mu$s trên \PCname{apple-m5}, $N = 1009$, $\mathrm{OW} = 0$: tỉ lệ
trên mọi lần trao lock (\%) / phần thời gian chạy nằm trong các lần đó (\%). Trung vị của 5 lần chạy.}
\label{tab:handoff}
\end{table}
% nguồn: scripts/report_numbers/extra_analyze.py (cột long% và time_in_long%).

Cùng giả thuyết phù hợp với việc MCS gãy mạnh hơn TAS và TTAS từ 4 lên 6 thread. MCS buộc lock đi qua
mọi thread, kể cả thread đang chạy trên E-core hay đang chờ core. TAS và TTAS cho phép thread vừa nhả
lấy lại lock. Ở 4 thread độ công bằng của chúng là 0,49 và 0,26, của MCS là 1,00.
% nguồn: scripts/report_numbers/rlock_analyze.py, fairness N=1009 OW=0 p=4: tas 0,49, ttas 0,26, mcs 1,00.

Lock thư viện cho thread chờ ngủ trong kernel, nên thread chờ không giành core của chủ lock và lock
không sụp. Thread đang ngủ cũng không đọc cache line của lock, nên thêm thread chờ gần như không thêm
chi phí. Một critical section của \PCname{std_mutex} tốn 1\,901~ns ở 8 thread và 2\,007~ns ở 20 thread.
% nguồn: scripts/report_numbers/rlock_analyze.py, ns/CS std_mutex N=1009 OW=0 p=8/20: 1\,901 / 2\,007.

Nhóm giả thuyết rằng mỗi lần đánh thức một thread qua kernel tốn một chi phí cố định cỡ 1,5~$\mu$s. Con
số này là ước lượng suy ra từ chính số đo lock. Báo cáo không có phép đo đánh thức riêng. Ở
$N = 1009$, mỗi critical section của \PCname{std_mutex} ở 20 thread tốn thêm 1\,459~ns so với khi không
tranh chấp (548~ns ở 1 thread, Bảng~\ref{tab:apple-m5:locksingle}, phụ lục). Ở $N = 10\,007$, mức tăng là 1\,469~ns, gần như không đổi dù phần việc dài gấp 10 lần. Mức
tăng đó bằng khoảng 0,28 lần một lời gọi \texttt{is\_prime(10\,007)}, khớp với mức 1,25--1,41 lần cận
dưới của lock thư viện từ 4 thread ở $N = 10\,007$. Nhóm cho rằng \PCname{pthread_mutex} (và \PCname{std_mutex} dựng
trên nó) trên macOS trao lock cho thread được đánh thức, nên gần như mỗi lần chuyển giao trả trọn một
lần đánh thức. Tới 10 thread, độ công bằng 0,88--0,97 của hai lock này phù hợp với điều đó. Khi vượt
số core, nó giảm dần xuống 0,78--0,79 ở 20 thread. Nhóm cho rằng khi đó scheduler chia core không đều
giữa các thread, độc lập với lock. Ở 2 thread mức tăng mới là 316~ns. Có thể thread chờ duy nhất nhiều lần chưa kịp ngủ
khi lock được nhả.
% nguồn: scripts/report_numbers/rlock_analyze.py, ns/CS std_mutex N=1009 OW=0 p=20: 2\,007 ns; locksingle 548; 2\,007 - 548 = 1\,459;
% p=2: 864,3 - 548,1 = 316. lock.csv (N=10007 OW=0, ns/CS): std_mutex p=1 5\,377, p=20 6\,847; 6\,847 - 5\,377 = 1\,469;
% 1\,469 / 5\,163 = 0,28. rlock_analyze.py độ chậm N=10007 OW=0 p=4..20: std_mutex 1,25..1,33, pthread_mutex 1,28..1,41 (1,41 ở p=16),
% os_unfair_lock 1,25..1,30. fairness N=1009 OW=0: std_mutex p=2..10 0,92..0,97, p=12/16/20 0,89/0,83/0,78;
% pthread_mutex p=2..10 0,88..0,96, p=12/16/20 0,87/0,86/0,79.

Tài liệu của \PCname{os_unfair_lock} (\texttt{os/lock.h}) ghi rằng thread vừa nhả có thể lấy lại lock
trước khi thread được đánh thức kịp thử. Phần lớn critical section vì thế không phải trả lần đánh thức
nào. Một critical section chỉ tốn thêm 302--351~ns từ 2 tới 20 thread, và độ công bằng giảm từ 0,97 ở 3
thread xuống 0,20 ở 20 thread.
% nguồn: MacOSX.sdk/usr/include/os/lock.h, đoạn @discussion của os_unfair_lock. scripts/report_numbers/rlock_analyze.py,
% ns/CS os_unfair_lock N=1009 OW=0: p=1 510,6; p=2..20 812,3..861,6 (nhỏ nhất p=2, lớn nhất p=16), trừ 510,6 = 301,7..351,0;
% fairness p=3 0,97, p=20 0,20.

Ở $N = 11$, \PCname{std_mutex} tốn thêm 19,7~ns mỗi critical section. Nếu mỗi lần đánh thức tốn
1\,459~ns như ở $N = 1009$, một lần đánh thức được chia cho khoảng 74 critical section. Nhóm cho rằng trong lúc thread được đánh thức chưa kịp chạy, thread vừa
nhả lock lấy lại lock và chạy tiếp các critical section của mình với cache line còn nóng. Chuỗi này kéo
dài khoảng 1,9~$\mu$s, cùng bậc với lần ngủ ngắn nhất của \texttt{sleep\_for} khi yêu cầu 100~ns
(2\,540~ns, Bảng~\ref{tab:apple-m5:sleeplatency}, phụ lục). Chi phí thêm mỗi critical section nhỏ hơn
một critical section của TAS khi tranh chấp (39,5~ns ở 4 thread, gồm 4,6~ns của
\texttt{is\_prime(11)}), nên phần lớn critical section không thể có chuyển giao. Độ công bằng 0,92 của
\PCname{std_mutex} ở 20 thread không chống lại giả thuyết. Thước đo này không thấy các đợt độc chiếm
luân phiên (mục Các yếu tố đe dọa độ tin cậy).
% nguồn: scripts/report_numbers/rlock_analyze.py, ns/CS std_mutex N=11 OW=0 p=20: 26,24; locksingle N=11 std_mutex 6,50;
% 26,24 - 6,50 = 19,74; 1\,459 / 19,74 = 74; 74 * 26,24 ns = 1,94 us. sleeplatency (results/apple-m5/sleep_latency.csv):
% sleep_for yêu cầu 100 ns, nhỏ nhất 2\,541 ns, trung vị 2\,667 ns. tas N=11 p=4: 39,53 ns/CS.
% fairness std_mutex N=11 OW=0 p=20: 0,92 (scripts/report_numbers/rlock_analyze2.py).

TAS thua TTAS nhiều nhất khi lock đổi chủ ở mỗi lượt. Phép nhả lock của TAS phải tranh cache line với
các \texttt{exchange} liên tục của thread chờ, còn thread TTAS chờ bằng phép đọc. Ở $\mathrm{OW} = 1$ và
8 thread, TAS đạt 0,467 triệu lần lấy lock mỗi giây, TTAS 1,41 triệu. Khi không có việc ngoài, lock chủ
yếu ở lại với thread vừa nhả. Độ công bằng của TAS 10 thread là 0,32 ở $\mathrm{OW} = 0$ và 0,91 ở
$\mathrm{OW} = 3$. Vì vậy TAS ở $\mathrm{OW} = 1$ và 3 chậm hơn ở $\mathrm{OW} = 0$ (6,79 và 7,17 so với
3,48 lần cận dưới ở 20 thread).
% nguồn: scripts/report_numbers/rlock_analyze.py, fairness N=1009 tas p=10: OW=0 0,32; OW=3 0,91. Thông lượng OW=1 p=8
% ttas 1,41 M/s, tas 0,467 M/s (locktime-ow1: 0,975 s và 2,94 s). Độ chậm tas p=20: OW=0 3,48, OW=1 6,79, OW=3 7,17.

Độ dài critical section quyết định tỉ lệ giữa chi phí chuyển giao và phần việc
(Hình~\ref{fig:apple-m5:lock_cs_length}, phụ lục). Ở $N = 11$ và 4 thread, một critical section của TAS tốn
39,5~ns, gấp khoảng 8,6 lần lời gọi \texttt{is\_prime(11)}. MCS chậm nhất ở điểm này (86~ns mỗi
critical section, 18,7 lần cận dưới). Nhóm cho rằng thứ tự FIFO buộc mọi critical section đổi chủ, còn
TAS và TTAS cho phép thread vừa nhả lấy lại lock thành từng đợt. Ở $N = 10\,007$, chi phí chuyển giao của
TAS chưa tới 1\,\% của 5,2~$\mu$s. Tới 4 thread mọi lock nằm trong 1,0--1,3 lần cận dưới. Từ 6 thread
TAS và TTAS tách dần (1,56 và 1,57 lần ở 10 thread). Khi vượt số core, khác biệt quay lại rõ. TAS và
TTAS chậm 2,5 lần ở 20 thread, lock thư viện 1,3 lần, còn MCS sụp. Nhóm cho rằng thread chờ quay giành core của
chủ lock, còn thread chờ của lock thư viện thì ngủ.
% nguồn: scripts/report_numbers/rlock_analyze.py ns/CS tas N=11 OW=0 p=4: 39,53; t(N=11) 4,60 ns (lockunit); 39,53/4,60 = 8,6;
% 39,53/5\,163 = 0,77 %. ns/CS mcs N=11 OW=0 p=4: 86,06; độ chậm 18,7. Độ chậm N=10007 OW=0 p=1..4: 1,00..1,32;
% p=10: tas 1,56, ttas 1,57; p=20: tas 2,55, ttas 2,53, std_mutex 1,33, pthread_mutex 1,28, os_unfair_lock 1,30, mcs 642.

Khi vượt số core, một critical section của MCS tốn gần như cùng thời gian ở $N = 1009$ và
$N = 10\,007$ (3,5 và 3,3~ms ở 20 thread), dù phần việc dài gấp 10 lần. Chi phí đó nằm ngoài phần việc.
Độ chậm của MCS ở $N = 10\,007$ (642 lần) thấp hơn ở $N = 1009$ (6\,880 lần) chỉ vì cận dưới lớn hơn.
% nguồn: lock.csv (mcs OW=0 p=20, ns/CS): N=1009 3,51e6, N=10007 3,31e6. rlock_analyze.py độ chậm mcs p=20 OW=0:
% N=1009 6\,880, N=10007 642.

\emph{Độ vững.} Cả bảy lần chạy MCS 10 thread đều chậm 7,1--15,5 lần cận dưới. Khi đo lại với load
average khoảng 3 thay vì khoảng 25, chỉ MCS ở 10 thread nhanh lên rõ, gần gấp đôi. Thời gian của các
lock khác bằng 0,83--1,05 lần thời gian ở đợt đo chính (mục Các yếu tố đe dọa độ tin cậy). Điều này phù hợp với cơ chế trên.
% nguồn: scripts/report_numbers/extra_analyze.py (lowload.csv): mcs p=10 low/main 0,53 (3\,545 / 6\,635 ns/CS);
% các dòng còn lại (kể cả mcs p<=8) 0,83 (tas 10t) .. 1,05 (std_mutex 4t).

Các điểm MCS vượt số core ở $N = 11$ không đo được sự sụp đổ, vì sau lần vượt thời hạn số vòng lặp giảm
tới mức hàng đợi không kịp hình thành. Ở $\mathrm{OW} \ge 1$, một số ô MCS vượt số core cho
kết quả hai cực, nên chỉ khoảng min--max đáng tin. Ở $\mathrm{OW} \ge 7$, trung vị ở 20 thread rơi vào cực
nhanh. Ở 16 thread, $\mathrm{OW} = 15$, bốn lần chạy xong dưới 10~ms, ba lần mất 2,0--10,6~s. Nhóm cho
rằng khi phần lớn thời gian nằm ngoài lock, hàng chờ thường ngắn. Lần chạy nào không có lượt trao lock
cho thread mất core thì xong nhanh, lần chạy nào có thì kéo dài nhiều giây.
% nguồn: lock.csv, mcs p=10 OW=0 N=1009 iters=34\,366: 1,241..2,727 s, cận dưới 0,1755 s.
% 6 dòng timeout: mcs N=11 OW=0/3 p=12/16/20, iters 4\,522/3\,391/2\,713; lần sau 45/33/27 vòng, ~0,1 ms.
% Hai cực (rlock_analyze.py, mcs N=1009, min..max s): OW=1 p=16 0,0071..7,96; OW=3 p=12 0,024..3,86, p=16 0,0057..8,75,
% p=20 0,0055..18,3; OW=7 p=20 trung vị 0,020 (0,0052..13,5); OW=15 p=20 trung vị 0,0075 (0,0066..17,4).
% mcs p=16 OW=15: 4 lần < 10 ms, 3 lần 2,0..10,6 s.

% ================================================================================================
\subsubsection{RQ3: Backoff giúp hay hại, và lock nhanh nhất có công bằng không?}

\begin{rqanswer}\textbf{Trả lời RQ3.} Backoff giúp TAS và gần như không đổi gì với TTAS. Ở 20 thread
($N = 1009$, $\mathrm{OW} = 0$), backoff bằng \texttt{sleep\_for} đưa TAS từ 12,2~s xuống 3,78~s
(thời gian quy về cùng số vòng lặp). Lợi ích đến từ việc lock gần như không đổi chủ, nên lock nhanh
nhất lại kém công bằng nhất. Khi có việc ngoài critical section và tranh chấp vừa phải, backoff làm
chậm. Hại lớn nhất là ở \PCname{ttas_bo_sleep} khi vượt số core với $\mathrm{OW} = 3$, chậm hơn TTAS
6,1--6,8 lần. Số liệu của \PCname{apple-m5}. RQ6 kiểm lại trên hai máy còn lại.\end{rqanswer}
% nguồn: locktime p=20: tas 12,2; tas_bo_sleep 3,78. rlock_analyze.py tỉ số ttas_bo_sleep/ttas N=1009 OW=3
% p=16/20: 6,13/6,78.

\PCfigure{apple-m5}{lock_backoff}

\emph{Bằng chứng.} Ở $\mathrm{OW} = 0$ (Hình~\ref{fig:apple-m5:lock_backoff}, cột trái),
\PCname{tas_bo_sleep} và \PCname{tas_bo_yield} chỉ mất 0,31--0,32 lần thời gian của TAS ở 20 thread,
còn \PCname{tas_bo_spin} mất 0,75 lần. Các biến thể của TTAS nằm trong 0,87--1,15 lần TTAS. Ở
$\mathrm{OW} = 3$ và 4--6 thread, backoff làm TAS chậm hơn (\PCname{tas_bo_sleep} 1,32--1,33 lần,
\PCname{tas_bo_spin} 1,29--1,49 lần). Ở 20 thread, \PCname{tas_bo_sleep} và \PCname{tas_bo_yield} lại
nhanh hơn (0,69 và 0,55 lần TAS), còn \PCname{tas_bo_spin} vẫn chậm hơn (1,18 lần). Trên
\PCname{LAP1184}, cũng ở $\mathrm{OW} = 3$, sleep backoff làm TAS chậm hơn cả ở số thread lớn (RQ6). Ở
$\mathrm{OW} = 15$, mọi biến thể nằm trong 0,90--1,26 lần TAS tới 12 thread. Ở 20 thread
\PCname{tas_bo_sleep} và \PCname{tas_bo_yield} nhanh hơn (0,73 và 0,63 lần).
% nguồn: locktime p=20: 3,78/12,2 = 0,31; 3,95/12,2 = 0,32. scripts/report_numbers/rlock_analyze.py tỉ số so với tas
% N=1009: OW=0 p=20 tas_bo_spin 0,745; tỉ số ttas_bo_*/ttas OW=0: 0,87..1,15. OW=3: tas_bo_sleep p=4/6 1,32/1,33,
% tas_bo_spin p=4/6 1,29/1,49; p=20 tas_bo_sleep 0,688, tas_bo_yield 0,549, tas_bo_spin 1,18.
% OW=15: p<=12 0,90..1,26; p=20 tas_bo_sleep 0,729, tas_bo_yield 0,633.
% LAP1184 (review: độ chậm N=1009 OW=3, rlock_analyze cho LAP1184): tas_bo_sleep/tas p=3..16 1,53..2,04, p=20 1,40, p=24 1,11.
% Ở OW=0 trên LAP1184 tas_bo_sleep vẫn nhanh hơn tas ở số thread lớn (1,21 so với 2,24 ở 24 thread), nên câu chỉ nói OW=3.

\PCfigure{apple-m5}{lock_fairness}

Ở $\mathrm{OW} = 0$, MCS giữ độ công bằng 1 tới 10 thread (Hình~\ref{fig:apple-m5:lock_fairness}). Độ
công bằng của \PCname{tas_bo_sleep} bằng đúng $1/p$, tới 0,05 ở 20 thread, tức là các thread chạy xong
lần lượt từng thread một. TAS và TTAS cũng thiên vị thread vừa nhả lock. TTAS thiên vị hơn TAS (0,26 so
với 0,49 ở 4 thread, 0,05 so với 0,11 ở 20 thread). Nhóm cho rằng phép nhả của TTAS không phải tranh
cache line với \texttt{exchange} nào, nên thread vừa nhả lấy lại lock trong lúc các thread chờ còn đang
đọc lại line. \PCname{os_unfair_lock}, lock thư viện nhanh nhất, kém công bằng nhất trong ba lock thư
viện (0,20 so với 0,78--0,79 ở 20 thread), vì lý do đã nêu ở RQ2.
% nguồn: scripts/report_numbers/rlock_analyze.py, fairness N=1009 OW=0: mcs p=2..10 1,00; tas_bo_sleep p=2/4/10/20 0,50/0,25/0,10/0,05;
% tas p=4/20 0,49/0,11; ttas p=4/20 0,26/0,05; os_unfair_lock p=20 0,20; std_mutex 0,78, pthread_mutex 0,79.

Ở $\mathrm{OW} = 3$, mọi lock gần như công bằng tới 4 thread. Từ 6 thread, $p$ vượt
$1+\mathrm{OW}$ nên luôn có thread chờ, và cùng sự thiên vị đó làm độ công bằng của TTAS giảm từ 0,98 ở 4 thread xuống
0,87 ở 6 thread và 0,55 ở 20 thread.
% nguồn: scripts/report_numbers/rlock_analyze.py, fairness N=1009 OW=3 ttas p=4/6/20: 0,98/0,87/0,55.

Khi $\mathrm{OW} > 0$, thread vừa nhả đi gọi \texttt{is\_prime} thay vì lấy lại lock ngay. Vì vậy
\PCname{os_unfair_lock} gần công bằng ở $\mathrm{OW} = 3$ và 15 (0,87--1,00), dù ở $\mathrm{OW} = 0$
nó chỉ còn 0,20 ở 20 thread.
% nguồn: lock.csv (N=1009, trung vị thr_min_s/thr_max_s), os_unfair_lock p=2..20: OW=3 0,87..1,00; OW=15 0,96..0,99;
% OW=7 0,94..1,00. Ngoại lệ OW=1: 0,99 ở 2 thread giảm còn 0,57 ở 20 thread (OW=1 không có trên hình).

\emph{Vì sao.} Các biến thể \PCname{tas_bo_*} lùi lại sau mọi lần \texttt{exchange} thất bại, kể cả khi
lock chỉ đang bị giữ, nên thread chờ gần như luôn ngủ hoặc nhường core. Các biến thể \PCname{ttas_bo_*}
chỉ lùi lại sau một lần \texttt{exchange} thua và vẫn đọc liên tục như TTAS (mục Cài đặt). Một lần
\texttt{sleep\_for} yêu cầu 100~ns thật ra ngủ trung vị 2\,670~ns. Ở $\mathrm{OW} = 0$, chủ lock nhả
rồi lấy lại lock ngay ở lượt kế tiếp của vòng lặp. Thread vừa thức dậy gần như luôn thấy lock bận và
ngủ tiếp với thời gian yêu cầu gấp đôi, tới khoảng 100~$\mu$s (\texttt{MAX\_DELAY} trong
\texttt{backoff.hpp}). Vì vậy một thread giữ lock tới hết vòng lặp của nó, rồi mới tới thread khác. Độ
công bằng đúng bằng $1/p$ phù hợp với điều này.
% nguồn: sleeplatency: sleep_for 100 ns trung vị 2\,667 ns. include/backoff.hpp: SleepBackoff bắt đầu
% PC_SLEEP_MIN_NS = 100, nhân đôi khi delay < PC_SLEEP_MAX_NS = 100000 ("MAX_DELAY of the assignment text").
% fairness tas_bo_sleep N=1009 OW=0 p=2..20: 0,50/0,34/0,25/0,16/0,12/0,10/0,08/0,06/0,05 (rlock_analyze.py).

\PCname{tas_bo_spin} chỉ giãn các lần thử. Thread chờ vẫn giữ core và vẫn gọi \texttt{exchange} sau
mỗi cửa sổ quay (tối đa 1\,024 lần \texttt{cpu\_relax()}, \texttt{backoff.hpp}). Nhóm cho rằng vì vậy
lợi ích của nó nhỏ hơn, và khi vượt số core nó còn chiếm core mà chủ lock cần.
% nguồn: include/backoff.hpp, SpinBackoff: cửa sổ ngẫu nhiên, nhân đôi tới PC_SPIN_MAX = 1024.

Khi $\mathrm{OW} > 0$, lock rảnh trong lúc thread chờ còn ngủ ít nhất 2,5~$\mu$s. Với
\PCname{tas_bo_spin}, cửa sổ quay lớn nhất kéo dài khoảng 10~$\mu$s (1\,024 lần \texttt{cpu\_relax()},
mỗi lần trung vị 9,6~ns), và lock cũng rảnh trong lúc đó. Ở tranh chấp vừa phải, đó là thời gian mất.
Ở số thread lớn, luôn có thread sẵn sàng lấy lock. Ở $\mathrm{OW} = 3$ và 15, khi vượt số core, thread
ngủ hoặc nhường core để lại core cho thread đang làm việc ngoài critical section. Nhóm cho rằng đó là lý
do \PCname{tas_bo_sleep} và \PCname{tas_bo_yield} nhanh hơn TAS ở 20 thread, còn \PCname{tas_bo_spin},
vẫn giữ core, thì không.
% nguồn: sleeplatency: sleep_for 100 ns nhỏ nhất 2\,541 ns; cpu_relax trung vị 9,6 ns; 1\,024 * 9,6 = 9,8 us.
% Tỉ số so với tas p=20: OW=3 tas_bo_sleep 0,688, tas_bo_yield 0,549, tas_bo_spin 1,18; OW=15 0,729 / 0,633 / 1,11.

MCS chia lượt đều khi mọi thread được chạy, và cả khi vượt số core nếu $\mathrm{OW} \le 3$ (độ công bằng
0,84--1,00). Khi vượt số core và $\mathrm{OW} \ge 7$, MCS kém công bằng nhất trong mười hai lock
(0,21--0,30 ở 16--20 thread, so với 0,78--0,81 của lock kế tiếp). Nhóm cho rằng một thread MCS bị lấy core giữa
\texttt{release} và \texttt{acquire} không nằm trong hàng đợi, nên các thread còn chạy lấy lock nhiều
lượt hơn. Khi $\mathrm{OW}$ lớn, khoảng giữa \texttt{release} và \texttt{acquire} dài, nên điều này xảy
ra thường hơn.
% nguồn: lock.csv (mcs N=1009, độ công bằng thr_min_s/thr_max_s): p<=10 mọi OW 0,84..1,00; p=12/16/20: OW=0 1,00/0,99/0,99,
% OW=1 1,00/0,96/0,95, OW=3 1,00/0,84/0,94; OW=7 p=16/20 0,30/0,21; OW=15 p=16/20 0,29/0,26 (rlock_analyze.py cùng số).
% Lock kế tiếp ở OW=7/15, p=16/20: 0,78..0,81.

Khi vượt số core và $\mathrm{OW} \ge 3$, \PCname{ttas_bo_sleep} chậm hơn TTAS 6,1--6,8 lần ở
$\mathrm{OW} = 3$ với 16--20 thread, 2,0--4,3 lần ở $\mathrm{OW} = 7$ và 1,6--1,8 lần ở
$\mathrm{OW} = 15$ (12--20 thread). Nhóm giả thuyết rằng ở biến thể này một phần thread chờ quay đọc và
chiếm core, phần khác ngủ. Thread vừa thức dậy sau khi ngủ có thể giành core của chủ lock. Trong lúc chủ
lock không chạy, các thread quay đọc dùng hết lượt của mình. \PCname{ttas_bo_yield}, chỉ khác ở chỗ
nhường core thay vì ngủ, không chậm hơn TTAS (1,97~s so với 2,42~s ở 20 thread, $\mathrm{OW} = 3$, thời
gian quy về cùng số vòng lặp). \PCname{tas_bo_sleep}, nơi mọi thread chờ đều ngủ, cũng không bị. Hai
số đo này phù hợp với giả thuyết.
% nguồn: scripts/report_numbers/rlock_analyze.py tỉ số ttas_bo_sleep/ttas N=1009: OW=3 p=16/20 6,13/6,78;
% OW=7 (độ chậm 16,6/3,89, 16,1/4,53, 11,5/5,85) p=12/16/20 4,27/3,55/1,97; OW=15 (5,63/3,28, 9,03/4,93, 9,1/5,82) 1,72/1,83/1,56.
% tables.tex locktime-ow3 p=20: ttas_bo_yield 1,97, ttas 2,42, ttas_bo_sleep 16,4; tas_bo_sleep
% 4,27 so với tas 6,20; p=16 ttas_bo_sleep 10,4, ttas 1,70.

\emph{Độ vững.} Kết quả của \PCname{ttas_bo_sleep} khi vượt số core phân tán mạnh giữa các lần lặp. Ở
16 thread và $\mathrm{OW} = 3$, thời gian quy về cùng số vòng lặp đi từ 1,0 tới 18,0~s, với trung vị
10,4~s. Đổi lệnh chờ \texttt{isb} thành \texttt{yield} không tạo
khác biệt vượt độ phân tán.
% nguồn: tables.tex locktime-ow3 p=16/20: ttas_bo_sleep 10,4/16,4, ttas 1,70/2,42; lock.csv ttas_bo_sleep
% p=16 OW=3: iters 21\,479, thô 0,259..4,510 s, quy về I = 85\,916 (x4,0): 1,04..18,04 s, trung vị 2,609 x 4 = 10,4 s.
% lock_yield.csv vs lock.csv: trung vị tỉ số 0,86..1,09, min..max 0,69..2,17.

\subsection{Barrier}
% TODO(report):
%  - Chi phí một episode tăng thế nào theo số thread (so với phân tích: một lần ghi phải lan tới
%    p - 1 core đang chờ).
%  - Biến thể spin sụp ở đâu khi vượt số core, và vì sao nặng hơn lock.
%  - Chờ chặn hoà vốn ở đâu; so với barrier thư viện.
%  - Backoff ở barrier có ý nghĩa gì khác so với ở lock.
% RQ4. Nguồn: results/apple-m5/barrier.csv (trung vị 7 lần), bảng barrier, sleeplatency.
\subsubsection{RQ4: Chi phí một barrier episode}

\begin{rqanswer}\textbf{Trả lời RQ4.} Với chờ spin, một episode tốn thêm khoảng 53~ns cho mỗi thread
trong cụm 4 P-core, dốc gần gấp 6 khi thread tràn sang E-core, và sụp khoảng 2\,000 lần khi số thread
vượt số core. Khi mỗi thread có core riêng, spin thuần nhanh nhất. Khi vượt số core,
\PCname{sr_bo_yield} nhanh nhất. Kết luận chỉ áp cho \PCname{apple-m5}.\end{rqanswer}

Nhóm đo số episode (một lần cả $p$ thread cùng qua barrier) mỗi giây khi giữa hai barrier không có việc
gì (Hình~\ref{fig:apple-m5:barrier_rate}, Bảng~\ref{tab:apple-m5:barrier}).

\PCfigure{apple-m5}{barrier_rate}
\PCtable{apple-m5}{barrier}

Với \PCname{sr}, mỗi thread thêm vào tốn khoảng 53~ns từ 2 tới 4 thread và khoảng 295~ns từ 4 tới
10 thread.
% nguồn: 1e9/trung vị ops_per_s của sr: p=2 94,5; p=4 200,1; p=10 1969 ns.
Phần tuyến tính đến từ $p$ phép \texttt{fetch\_add} lần lượt trên cache line của bộ đếm và $p-1$ lần
đọc lại cờ sau khi cờ đảo. Chỗ gãy sau 4 thread trùng với số P-core, và kết quả lock gãy ở cùng chỗ
ngay cả ở ô đối chứng mà lock gần như rảnh (RQ2). Vì vậy nó đến từ máy chứ không từ thuật toán barrier.
Nhóm giả thuyết hai nguyên nhân là E-core chậm hơn và cache line phải đi giữa hai cụm L2. Vì không gắn
được thread vào core, nhóm không tách được hai nguyên nhân này.

Từ 10 lên 12 thread, \PCname{sr} chậm đi khoảng 2\,070 lần, và mỗi thread vượt số core thêm khoảng
2~ms vào một episode.
% nguồn: sr: tốc độ p=10/p=12 = 2067; ms mỗi episode/(p-10), p=12..20: 1,7-2,0.
Một episode cần \emph{mọi} thread được chạy, nhưng thread đã tới quay giữ core tới hết lượt lập lịch.
Barrier giống MCS ở điểm này (RQ2).

Chờ chặn hoà vốn đúng tại số core. Dưới đó, \PCname{sr_block} luôn chậm hơn \PCname{sr}, vì đánh thức
qua kernel tốn vài micro giây. Ở 12 thread, nó nhanh hơn khoảng 217 lần.
% nguồn: trung vị ops_per_s p=12: sr_block 53319, sr 245,7.
\PCname{sr_bo_yield} nhanh hơn \PCname{std_barrier} ở mọi số thread từ 2 tới 20, và nhanh nhất trong
mọi barrier khi vượt số core (11,6~$\mu$s mỗi episode ở 12 thread). Thread chờ gọi
\texttt{sched\_yield}, nhường core cho đúng thread chưa tới barrier. Nhờ vậy episode xong mà không phải chờ hết
lượt lập lịch, và không phải trả một lần đánh thức qua kernel như \PCname{sr_block}. Khi mỗi thread có
core riêng, lời gọi hệ thống này chỉ làm thread chậm nhận ra cờ đã đảo, nên \PCname{sr_bo_yield} chậm hơn
\PCname{sr} tới 1,7 lần (8 thread).
% nguồn: ns/episode (1e9/trung vị ops_per_s): p=12 sr_bo_yield 11555, std_barrier 13657, sr_block 18755; p=8 sr_bo_yield 1225, sr 733.

Backoff làm barrier chậm đi khi mỗi thread có core riêng (\PCname{sr_bo_spin} chậm hơn \PCname{sr}
2,3 tới 9,5 lần), vì thread chờ chỉ đọc cờ nên không có tranh chấp để bớt. Một episode của
\PCname{sr_bo_sleep} ở 2 thread tốn 2,49~$\mu$s, xấp xỉ thời gian ngủ thật ngắn nhất của
\texttt{sleep\_for} khi yêu cầu 100~ns (2,54~$\mu$s).
% nguồn: sr_bo_sleep p=2: 1e9/401186 = 2493 ns; sleep_for 100 ns: nhỏ nhất 2540 ns.

Đối chiếu với dự đoán, D4.1 chỉ đúng từng đoạn, vì chi phí tăng tuyến tính theo $p$ nhưng độ dốc đổi
ở 4 thread. D4.2 đúng. Spin sụp ngay khi vượt số core, và chờ chặn chỉ thắng ở vùng đó. D4.3 đúng.
Backoff làm barrier chậm khi mỗi thread có core riêng. Các ô sụp đổ dựa trên 50--69 episode mỗi lần chạy nên chỉ cho bậc độ lớn. Khoảng
cách 2--3 bậc vẫn lớn hơn nhiều độ phân tán giữa 7 lần lặp.
% nguồn: cột iters của sr, p=12..20.

\subsection{Concurrent counter}
% TODO(report):
%  - Thứ tự nhanh chậm giữa fetch-and-add, CAS và counter dùng lock, ở tranh chấp cao và thấp.
%  - Số lần CAS hỏng mỗi lần tăng: bằng chứng trực tiếp vì sao CAS chậm hơn fetch-and-add.
%  - Backoff giúp hay hại; seq_cst có khác bản mặc định không (và vì sao, theo lệnh máy).
%  - Counter dùng lock: liên hệ với kết quả của phần lock ở critical section rất ngắn.
% RQ5. Nguồn: results/apple-m5/counter.csv (trung vị 7 lần, ok=1).
\subsubsection{RQ5: Thứ tự hiệu năng của các counter}

\begin{rqanswer}\textbf{Trả lời RQ5.} Khi không có việc giữa hai lần tăng, fetch-and-add nhanh hơn CAS
2,1 tới 2,7 lần từ 4 thread, vì CAS hỏng và phải thử lại. Counter dùng spinlock không backoff chậm
nhất. Mutex thư viện và CAS có backoff lại nhanh hơn fetch-and-add ở chế độ này vì chúng dồn các lần
tăng về một core. Lợi thế này mất khi có việc giữa hai lần tăng. Kết luận chỉ áp cho \PCname{apple-m5}.\end{rqanswer}

Nhóm đo thời gian để mỗi thread tăng counter $10^6$ lần, với $\mathrm{OW}$ lời gọi
\texttt{is\_prime(101)} giữa hai lần tăng (Hình~\ref{fig:apple-m5:counter_time},
Bảng~\ref{tab:apple-m5:countertime}). Nhóm cũng đo thời gian trung bình một lần tăng và độ công bằng
như ở RQ3.

\PCfigure{apple-m5}{counter_time}

Không counter nào mở rộng. Với $\mathrm{OW} = 0$, thông lượng của \PCname{faa} giảm từ khoảng 610 triệu
lần tăng mỗi giây với 1 thread xuống 42 triệu với 20 thread. Nguyên nhân là mỗi lần tăng phải giữ
cache line của counter ở trạng thái độc quyền.
% nguồn: faa OW=0, seconds/ops: p=1 1,64 ns; p=20 23,6 ns.

\PCtable{apple-m5}{countertime}

\PCname{cas} chậm hơn \PCname{faa} 2,1 tới 2,7 lần từ 4 thread. Từ 8 thread, số lần CAS hỏng trên mỗi
lần tăng vượt 1 (Hình~\ref{fig:apple-m5:counter_cas_failures}).
% nguồn: t(cas)/t(faa) OW=0 p=4..20; cas_fail/ops p=8 1,59.

\PCfigure{apple-m5}{counter_cas_failures}

Một lần CAS hỏng đã tốn một lần kéo cache line về mà không tăng được counter. Lệnh \texttt{ldadd} của
fetch-and-add thì đọc và ghi trong một bước. Nếu mỗi lần hỏng tốn một lần chuyển, thời gian của \PCname{cas} xấp xỉ thời gian của \PCname{faa} nhân
$1+f$, với $f$ là số lần hỏng mỗi lần tăng. Từ 6 tới 20 thread, mô hình lệch tỉ số đo được không quá 15\%.

Ở 2 thread, ngược với D5.1, \PCname{cas} nhanh hơn \PCname{faa} khoảng 28\%. Ở đó chỉ có 0,09 lần
CAS hỏng trên mỗi lần tăng, nên gần như mọi lần tăng thành công ngay lần đầu. Mô hình $1+f$ không giải
thích được phần chênh còn lại. Kiến trúc ARMv8.1 cho phép phần cứng thực hiện \texttt{ldadd} ở gần cache
dùng chung thay vì trong L1 của core gọi. Nếu M5 làm vậy, mỗi lần \texttt{ldadd} tốn một lượt đi về tới
cache dùng chung. Một thread chạy vòng CAS thì có thể tăng vài lần liên tiếp trên line đang nằm sẵn trong
L1 của nó. Đây là giả thuyết. Ghi lại giá trị trả về của
từng lần tăng sẽ cho biết mỗi thread tăng liên tiếp bao nhiêu lần trước khi line rời core của nó. Con
số đó cho phép kiểm giả thuyết này.
% nguồn: t(cas)/t(faa) OW=0 p=2: 4,21/5,85 = 0,72; cas_fail/ops p=2: 0,091.

Khi $\mathrm{OW} \le 1$, counter dùng spinlock không backoff chậm nhất. Ở 20 thread, \PCname{lock_tas}
chậm hơn \PCname{faa} 13,7 lần. \PCname{lock_mcs} chậm hơn \PCname{lock_tas} 3,2 lần ngay ở 2 thread.
MCS trao lock theo thứ tự đến, nên mỗi lần tăng chuyển lock và giá trị sang core khác.
% nguồn: OW=0: lock_tas/faa p=20 13,67; lock_mcs/lock_tas p=2 3,23.

Ngược với D5.1, ở $\mathrm{OW} = 0$ counter dùng \PCname{lock_pthread_mutex} và
\PCname{lock_std_mutex} nhanh hơn \PCname{faa} từ 4 thread. Ở 8 thread, một lần tăng tốn 12,4 và 14,1~ns,
so với 23,4~ns của \PCname{faa}. Với \PCname{faa}, mọi thread tiến đều nhau nên gần như mỗi lần tăng phải
kéo line của counter từ core khác về. Một lần tăng dưới mutex rẻ hơn thế, nên tính trung bình nó tốn ít
hơn một lần chuyển line. Người giữ lock thực hiện nhiều lần tăng liên tiếp trên line đã nằm trong cache
của mình. Các thread chờ ngủ trong kernel và không tranh line. Độ công bằng theo thời điểm kết
thúc của mutex vẫn khoảng 0,98, nên các chuỗi lần tăng liên tiếp này ngắn so với cả lần chạy. Khi có
việc giữa hai lần tăng ($\mathrm{OW} = 8$, 20 thread), mỗi lần lấy lock gần như luôn phải đánh thức một
thread đang ngủ. Khi đó \PCname{lock_pthread_mutex} chậm hơn \PCname{faa} 4,9 lần.
% nguồn: OW=0 p=8 ns/op: pthread_mutex 12,41, std_mutex 14,10, faa 23,39; thr_min/thr_max mutex 0,97-0,98, faa 0,94; OW=8 p=20: 9,98 s / 2,03 s.

Backoff dồn các lần tăng về một thread tại mỗi thời điểm. Thời gian một lần tăng của
\PCname{cas_sleep} bằng giá trị 1 thread trong khoảng 11\% ở mọi số thread và mọi $\mathrm{OW}$. Độ
công bằng của nó xấp xỉ $1/p$, đúng như khi các thread chạy lần lượt. Vì vậy, ở 20 thread,
\PCname{cas_sleep} nhanh hơn \PCname{cas} 32 lần khi $\mathrm{OW} = 0$ nhưng chậm hơn 3,1 lần khi
$\mathrm{OW} = 8$. Trong trường hợp sau, chạy tuần tự bỏ phí phần \texttt{is\_prime} lẽ ra chạy song song.
% nguồn: cas_sleep seconds/ops lệch p=1 tối đa 1,107; thr_min/thr_max p=20 0,049.

Thứ tự \texttt{seq\_cst} không làm chậm fetch-and-add, nhưng làm \PCname{cas} chậm khoảng 2 lần ở
$\mathrm{OW} = 0$ và gần gấp đôi số lần hỏng. Ở $\mathrm{OW} \ge 1$, khác biệt biến mất. Nhóm giả thuyết rằng
\texttt{ldar} và \texttt{casal} ngăn các lần lặp chồng lên nhau.
% nguồn: t(cas_seqcst)/t(cas) OW=0: 1,29-2,61; OW=1,8: 0,83-1,10; t(faa_seqcst)/t(faa): 0,78-1,23.

Đối chiếu với dự đoán, D5.1 đúng với \PCname{faa}, \PCname{cas} và spinlock từ 4 thread, nhưng sai ở hai
điểm đã nêu (mutex thư viện ở $\mathrm{OW} = 0$, và \PCname{cas} ở 2 thread). D5.2 đúng. Số lần CAS hỏng
trên mỗi lần tăng tăng từ 0,09 ở 2 thread lên 1,59 ở 8 thread. D5.3 đúng và mạnh hơn dự đoán. Backoff
đưa số lần hỏng về gần 0, nhưng bằng cách cho các thread chạy lần lượt, nên mất công bằng. D5.4 sai với
CAS, vì \texttt{seq\_cst} làm CAS chậm khoảng 2 lần khi không có việc giữa hai lần tăng.

Nhóm không phân tích \PCname{lock_mcs} trên 10 thread, vì các ô đó dựa trên 20--200 lần tăng mỗi thread và có
3 lần vượt thời hạn (Bảng~\ref{tab:apple-m5:campaign}).

\subsection{So sánh giữa các máy}
\PCfigureopt{cross}{lock}
\PCfigureopt{cross}{barrier}
\PCfigureopt{cross}{counter}
% TODO(report): kết quả có giống nhau giữa các máy không; nếu khác thì khác vì kiến trúc (x86 so
% với ARM), số core, P/E-core hay SMT, hay hệ điều hành (scheduler, sleep, ghim thread). Mục này
% chỉ có hình khi results/ chứa từ hai máy trở lên.
\subsubsection{RQ6: Kết luận trên máy khác}

\begin{rqanswer}\textbf{Trả lời RQ6.} Chưa kiểm được, vì báo cáo mới có số đo của \PCname{apple-m5}. Hình so
sánh tự xuất hiện khi có từ hai máy (Bảng~\ref{tab:cross:machines}).\end{rqanswer}

Một máy x86 kiểm được ba giả thuyết. Với TSO, mọi phép đọc-sửa-ghi atomic là lệnh có tiền tố
\texttt{lock}, nên \PCname{cas_seqcst} phải bằng \PCname{cas}. Nếu không, giải thích ở RQ5 sai. Với SMT,
chi phí barrier và counter phải gãy ở số core vật lý, còn spin barrier chỉ sụp khi vượt số luồng phần
cứng. Độ trễ của \texttt{pause} và thời gian ngủ thật trên Linux đổi chi phí của các biến thể backoff.

\TODO{ghi kết quả từng giả thuyết trên máy thứ hai và các khác biệt lớn nhất giữa các máy.}

% ====================================================================================================
\section{Các yếu tố đe dọa độ tin cậy}
% TODO(report): mỗi yếu tố: nó có thể làm sai kết quả thế nào, nhóm đã làm gì để giảm, còn lại gì.
%  - Tải nền và nhiệt độ (throttling); load average lúc bắt đầu và kết thúc đợt đo (campaign.log).
%  - Hai loại core (P/E) và không ghim được thread trên macOS: cận dưới dùng t đo trên một core nhanh.
%  - Cận dưới lý tưởng chứa chi phí lấy/nhả lock khi không tranh chấp, nên hơi rộng.
%  - Số vòng lặp khác nhau giữa các cấu hình và phép quy đổi tuyến tính; cấu hình đo với I giảm.
%  - Lần chạy vượt thời hạn: chỉ biết cận dưới của thời gian.
%  - Số lần lặp và độ phân tán; trung vị che gì.
%  - Chỉ đo một compiler và một bộ cờ trên mỗi máy; máy ảo hoặc container (nếu dùng).
%  - Giới hạn của kiểm tra đúng đắn (mục 4).
%  - Tính đại diện của workload is_prime.
% Mục 7: Các yếu tố đe dọa độ tin cậy, ba đoạn Construct/Internal/External. Mỗi giới hạn chỉ nêu ở
% đây. Nguồn: results/{apple-m5,C4M7Y670NT,LAP1184}/{campaign.log, machine.txt, lock.csv, counter.csv,
% barrier.csv, calibration.json, tsan.log}, tables.tex (campaign, lockunit, locksingle), results/apple-m5/extra/.
% Số tự tính: scripts/report_numbers/corr_stats.py, corr_stats2.py, extra_analyze.py, rlock_analyze.py.

\emph{Construct validity.} Cận dưới của phần lock dùng $t$ đo ở 1 thread, nên chứa chi phí lấy và nhả
lock của lock nhanh nhất. Ở $N=11$, chi phí này chênh 41\,\% giữa các lock
(Bảng~\ref{tab:apple-m5:locksingle}, phụ lục). Cận dưới bỏ qua giới hạn 10 core khi $p > 10$ và
$\mathrm{OW} = 15$ (ở 20 thread sàn thật cao hơn 1,6 lần). Cận dưới còn giả định mọi core chạy
\texttt{is\_prime} nhanh như ở lần chạy 1 thread. Khi số thread vượt số P-core, một phần việc chạy trên
E-core chậm hơn, nên một phần độ chậm không do lock. Ở 6 thread và $\mathrm{OW} = 15$, khi lock chưa
bận liên tục, mọi lock đã chậm 1,53--2,03 lần cận dưới.
% nguồn: locksingle N=11: 6,50/4,60 = 1,41; sàn khi p > 10: p(1+OW)It/10 so với max(p,1+OW)It, chỉ lớn hơn
% khi OW = 15 (p=12: 19,2/16 = 1,2; p=16, 20: 25,6/16 = 32/20 = 1,6); OW <= 7 thì p(1+OW)/10 <= p.
% scripts/report_numbers/rlock_analyze.py, độ chậm N=1009 OW=15 p=6: 1,53 (mcs) .. 2,03 (os_unfair_lock).
Phép quy về cùng số vòng lặp giả định thời gian tỉ lệ với số vòng. Giả định này yếu nhất ở MCS quá tải
(18--286 vòng mỗi thread trên \PCname{apple-m5}, 15--195 trên \PCname{C4M7Y670NT}, 10--56 trên
\PCname{LAP1184}), ở counter \PCname{lock_mcs} quá tải (10--200 lần tăng mỗi thread trên cả ba máy) và
ở các ô barrier sụp (50--69 episode mỗi lần chạy). Trên \PCname{LAP1184}, ở $N=1009$ và
$\mathrm{OW}=0$, phần lớn lần chạy MCS ở 16 và 24 thread xong trong vài mili giây, còn số ít mất 1,2--2,6~s.
Trung vị của hai ô này rơi vào nhóm nhanh, nên chỉ là cận dưới của chi phí MCS. Vài ô MCS quá tải chỉ
còn 4--6 giá trị sau khi vượt thời hạn (mục Phương pháp đo).
% nguồn: plot.aggregate(peers=True) trên lock.csv, mcs reduced: iters apple-m5 18..286, C4M7Y670NT 15..195,
% LAP1184 10..56 (N=10007: 10); counter.csv lock_mcs iters <= 200: apple-m5 từ 12 thread, C4M7Y670NT từ
% 16 thread, LAP1184 từ 10 thread, iters 10..200; barrier.csv (sr, sr_bo_spin, 12-20 thread: iters 69/66/50);
% LAP1184 lock.csv mcs N=1009 OW=0: p=16 5 lần 2,7..6,2 ms, 2 lần 1,18/1,73 s; p=24 6 lần 3,6..5,5 ms,
% 1 lần 2,58 s.
Ở $\mathrm{OW}=8$, thời gian của mọi counter, kể cả \PCname{faa}, tăng ít nhất gấp đôi từ 4 lên 10 thread,
và nhóm không có phép đo đối chứng để tách phần do E-core. Phép đo khoảng trống trao lock thêm hai lần
đọc đồng hồ vào mỗi critical section, nên báo cáo chỉ dùng thống kê khoảng trống của nó, không dùng thời
gian chạy. Độ công bằng so thời điểm kết thúc của các thread, nên không thấy các đợt độc chiếm luân phiên
ngắn.
% nguồn: counter.csv OW=8, giây/iters trung vị 10 thread / 4 thread: faa 2,13 (nhỏ nhất) .. lock_mcs 6,44;
% experiments/handoff_gap.cpp; định nghĩa độ công bằng: bench_lock.cpp, bench_counter.cpp (thr_min_s, thr_max_s).

\emph{Internal validity.} Không lần chạy nào ghim thread vào core. \PCname{apple-m5} có 4 P-core và 6
E-core, và macOS không cho ghim thread. Ở $N=1009$, thời gian một lời gọi \texttt{is\_prime} trong 420
lần chạy 1 thread đi từ 468 tới 1\,713~ns. Nhóm cho rằng độ phân tán này phần lớn do thread rơi vào
E-core, nhưng không tách được ảnh hưởng của E-core với tải nền. Load average 1 phút trong đợt đo lock
có trung vị khoảng 25 ở mọi số thread, kể cả lần chạy 1 thread (12,4--51,2). Con số này cao gấp đôi số
thread trung bình theo thời gian của đợt đo lock (12,0).
% nguồn: lock.csv cột load1: 12,42..51,18, trung vị 25,4 (p=1: 25,2) (corr_stats.py); số thread trung bình
% có trọng số seconds trên mọi dòng lock.csv apple-m5: 11,97 (chỉ dòng ok: 11,78).
Nhóm đo lại năm lock (\PCname{tas}, \PCname{ttas}, \PCname{mcs}, \PCname{std_mutex},
\PCname{os_unfair_lock}) ở $N = 1009$, $\mathrm{OW} = 0$, 1--10 thread, 7 lần lặp, mỗi lần chạy chỉ bắt
đầu khi load average dưới 3 hoặc sau 2 phút chờ (trung vị 2,9). Ở cấu hình này, các đường cong giữ hình
dạng khi tải thấp, còn điểm gãy của MCS ở 10 thread nhỏ đi một nửa (0,53 lần, tức chậm khoảng 7 lần
cận dưới thay vì 13). Ngoài MCS ở 10 thread, trung vị mới bằng 0,83--1,05 lần trung vị của đợt đo
chính, tức lệch tới 17\,\%. Báo cáo không diễn giải những khác biệt nhỏ hơn mức này. Lần đo lại khác
đợt đo chính cả về load average lẫn phiên chạy và khoảng nghỉ giữa các lần chạy, nên không tách được
riêng ảnh hưởng của tải. Vùng vượt số core chưa được đo lại.
% nguồn: scripts/report_numbers/extra_analyze.py (results/apple-m5/extra/lowload.csv, 175 lần chạy, load1
% 2,77..6,71, trung vị 2,89); ns/CS low/main ngoài mcs p=10: 0,83 (tas 10t) .. 1,05 (std_mutex 4t);
% mcs p=10: 3\,545 / 6\,635 ns = 0,53; t(N=1009) = 6\,635/13,0 = 510 ns -> 3\,545/510 = 6,9.
Nhóm không đọc nhiệt độ hay tần số CPU. Thời gian 1 thread không trôi quá 3,4\,\% qua 7 lượt lặp. Với 7
lần lặp, nhóm không tính khoảng tin cậy. Trung vị không thể hiện phân bố hai đỉnh, và chỉ thanh min--max
cho thấy phần này.
% nguồn: scripts/report_numbers/corr_stats2.py (lock.csv, threads=1, n=1009, giây/(iters*(1+OW)));
% scripts/report_numbers/corr_stats.py (tỉ số theo lượt 0,977..1,034)
Trên \PCname{LAP1184}, load average (5,90--19,6) chỉ đếm tiến trình trong máy ảo. Tải của Windows, việc
Windows xếp lịch cho các CPU ảo, nguồn điện và chế độ điện trong lúc đo không được ghi lại. Linux không
biết CPU ảo nào nằm trên P-core hay E-core, và \texttt{lscpu} báo tần số cố định 2\,496~MHz, nên không
thấy được throttling. Thời gian một lời gọi ở 1 thread, $N=1009$, đi từ 1\,463 tới 3\,176~ns qua 455
lần chạy. \texttt{sleep\_for} với yêu cầu 100~ns ngủ trung vị khoảng 62~\textmu s. Trong con số này,
nhóm không tách được phần của bộ hẹn giờ ảo của Hyper-V với phần của timer slack, tức khoảng trễ Linux
được phép thêm vào khi đánh thức để gộp các lần hẹn giờ. Pause-loop exiting, cơ chế hypervisor tạm dừng
một CPU ảo quay vòng chờ quá lâu, và việc đánh thức thread qua ngắt ảo là những giải thích khác, chưa
đo, cho việc barrier spin chậm ngay ở 12 thread.
% nguồn: figures/LAP1184/tables.tex (campaign); machine.txt (cpu MHz 2496.0, virtualization: Microsoft);
% lock.csv LAP1184 threads=1, n=1009: 1463..3176 ns/lời gọi, 455 lần chạy; sleep_latency.csv
% (sleep_for requested_ns=100: median_ns 61655).

\emph{External validity.} Số đo đến từ ba máy, gồm hai máy Apple Silicon (Apple clang 21, macOS) và một máy
x86-64 (g++ 11.5, Linux trong máy ảo WSL2), cả ba đều \texttt{-O2}. Hai máy Apple khác nhau ở số P-core
và số core, nhưng cũng ở thế hệ chip (M3 Pro, M5), dung lượng L2 mỗi cụm và phiên bản macOS. Vì vậy việc
điểm gãy dời theo số P-core phù hợp với giải thích hai loại core, nhưng không chứng minh nó. Máy x86-64
khác hai máy Apple cùng lúc ở kiến trúc, memory model, compiler, hệ điều hành và việc có máy ảo hay
không. Vì vậy RQ6 chỉ gán một khác biệt của máy này cho một nguyên nhân khi có bằng chứng riêng, như
\texttt{seq\_cst} với TSO. Nhóm không có máy x86 chạy Linux trực tiếp, máy chạy Windows trực tiếp, hay
máy nhiều socket. \texttt{is\_prime} không chạm dữ liệu dùng chung. Critical section thật thường ghi dữ
liệu dùng chung, nên chi phí chuyển giao lock thực tế có thể lớn hơn.
% nguồn: machine.txt apple-m5 (l2_cache 16 MiB per 4 P cores / 6 MiB per 6 E cores; macOS 27.0.1) và
% C4M7Y670NT (16 MiB per 6 P cores / 4 MiB per 6 E cores; macOS 26.7.1).

% ====================================================================================================
\section{Kết luận}
% TODO(report): trả lời từng câu hỏi nghiên cứu của mục 1 bằng một hai câu có số; điều học được
% khuyến nghị dùng lock, barrier, counter nào trong hoàn cảnh nào; việc chưa làm.
% Kết luận. Mọi con số dưới đây đã xuất hiện ở mục Kiểm tra tính đúng đắn hoặc mục Kết quả và phân tích.

Nhóm cài đặt ba spinlock, một sense-reversing barrier và ba họ concurrent counter, kiểm tra chúng bằng
stress test, ThreadSanitizer và mutation testing, rồi đo chúng trên \PCname{apple-m5} và
\PCname{LAP1184}. Câu trả lời cho RQ1 tới RQ5 dựa trên \PCname{apple-m5}. RQ6 so hai máy.

RQ1. Mọi cài đặt qua mọi bài kiểm tra, và bộ kiểm tra bắt được cả 14 biến thể cố ý làm hỏng. Không kỹ
thuật nào tự nó đủ. Lỗi thuật toán lộ ra dưới stress test, còn lỗi memory order như nhả lock bằng store
\texttt{relaxed} chỉ ThreadSanitizer bắt được.

RQ2. Tới 4 thread, cả ba spinlock bám cận dưới lý tưởng trong khoảng 1,24 lần. Ở 20 thread, TTAS chậm
2,48 lần cận dưới và TAS 3,48 lần, đúng thứ tự mà lưu lượng coherence dự đoán. MCS sụp ngay khi số thread
bằng số core, vì nó trao lock theo thứ tự cố định. Mỗi lần thread kế tiếp không có core là một lượt lập
lịch bị mất. Lock thư viện không sụp vì thread chờ ngủ trong kernel.

RQ3. Backoff giúp TAS nhiều nhất. Ở 20 thread, TAS với sleep backoff chỉ tốn 0,31 lần thời gian của TAS.
Nhưng nó đạt điều đó bằng cách cho các thread chạy lần lượt, nên độ công bằng đúng bằng $1/p$. Lock nhanh
nhất là lock kém công bằng nhất. MCS công bằng tuyệt đối chừng nào mỗi thread còn có core.

RQ4. Chi phí một episode của barrier tăng khoảng 53~ns cho mỗi thread trong cụm P-core, dốc hơn khi
thread tràn sang E-core, và tăng khoảng 2\,000 lần khi số thread vượt số core. Spin thuần là lựa chọn tốt
khi mỗi thread có core riêng. Khi vượt số core, chờ bằng \texttt{sched\_yield} nhanh nhất.

RQ5. Không counter nào mở rộng. Fetch-and-add nhanh hơn vòng CAS 2,1 tới 2,7 lần từ 4 thread vì số lần
CAS hỏng trên mỗi lần tăng tăng từ 0,09 ở 2 thread lên 1,59 ở 8 thread. Backoff đưa số lần hỏng về gần 0
nhưng tuần tự hoá các thread. Mutex thư viện thắng fetch-and-add khi không có việc giữa hai lần tăng,
vì người giữ lock tăng nhiều lần liên tiếp trên cache line nằm sẵn trong cache của mình.

RQ6. Một phần. Trên \PCname{LAP1184} (x86-64, Linux trong máy ảo WSL2), mọi cài đặt cũng qua mọi bài
kiểm tra, và các thứ hạng chính vẫn giữ: TTAS nhanh hơn TAS, MCS sụp khi mọi luồng phần cứng đều bận,
fetch-and-add nhanh hơn vòng CAS. Ba khác biệt đến từ nền tảng. Với TSO, \texttt{seq\_cst} không tốn
thêm gì (0,75 tới 1,14 lần \PCname{cas}), và stress test không bao giờ bắt được lỗi memory order. Barrier
spin sụp ngay ở 12 thread, bằng số luồng phần cứng, có lẽ vì Windows phải lấy core của các CPU ảo.
\texttt{sleep\_for} ngắn nhất kéo dài khoảng 61~\textmu s, gấp 23 lần trên \PCname{apple-m5}, nên các
biến thể ngủ chậm hơn nhiều.

Bảy dự đoán viết trước khi đo bị số liệu bác bỏ toàn phần hoặc một phần. MCS không phải spinlock nhanh
nhất khi số thread chưa vượt số core (D2.1). Lock chỉ bão hoà tại $p = 1 + \mathrm{OW}$ khi
$1 + \mathrm{OW} \le 4$ (D2.3), và sleep backoff vẫn giúp TAS ở nhiều thread dù có việc ngoài
critical section (D3.2). Chi phí barrier không tăng theo một độ dốc duy nhất (D4.1).
Mutex thư viện và vòng CAS ở 2 thread nhanh hơn fetch-and-add (D5.1). \texttt{seq\_cst} làm CAS chậm
khoảng 2 lần (D5.4). Khi vượt số core, MCS mất cả tính công bằng (D3.3). Điểm chung của các dự đoán sai
là chúng chỉ tính lưu lượng coherence mà bỏ qua scheduler của hệ điều hành và việc máy có hai loại core.

Từ các kết quả này, nhóm khuyến nghị ba điều. Khi số thread có thể vượt số core, dùng lock thư viện
hoặc TTAS có backoff, và chỉ dùng MCS khi mỗi thread chắc chắn có core riêng. Dùng barrier spin khi số
thread không vượt số core, và barrier có \texttt{sched\_yield} khi có thể vượt. Dùng fetch-and-add cho
counter, và chia counter theo thread khi nó nằm trên đường nóng. Việc chưa làm gồm phiên bản phân tán bằng MPI-3 RMA
(phần 4 của đề), số đo trên các máy còn lại của nhóm, và một máy x86 chạy Linux trực tiếp để tách ảnh
hưởng của máy ảo.

% ====================================================================================================
% TODO(report): thay \nocite{*} bằng các \cite đặt đúng chỗ trong bài (thuật toán và ý tưởng lấy
% từ sách, bài báo, bài giảng phải được trích nguồn); xoá mục nào trong references.bib không dùng.
\nocite{*}
\clearpage
\phantomsection\addcontentsline{toc}{section}{References}
\bibliographystyle{plain}
\bibliography{references}

% ====================================================================================================
\appendix
\section{Cách build và chạy lại}
Mã nguồn cần compiler C++20 (Apple clang, g++ hoặc clang++), \texttt{make} và Python~3. Phần vẽ hình
cần thêm \texttt{matplotlib} và \texttt{numpy}. Báo cáo biên dịch bằng XeLaTeX (TeX Live). Chi tiết
trong \texttt{README.md}.

\begin{lstlisting}[language=bash]
make                 # biên dịch các bài kiểm tra và các benchmark vào build/
make test            # chạy mọi bài kiểm tra tính đúng đắn
make tsan            # chạy lại dưới ThreadSanitizer
make mutation        # các biến thể cố ý làm hỏng phải bị bộ kiểm tra bắt
make check           # bốn lệnh trên, theo thứ tự
make bench          # chạy toàn bộ đợt đo trên máy rảnh -> results/<máy>/
make figures         # vẽ lại mọi hình và bảng từ results/ -> report/figures/
make report          # biên dịch report/report.pdf
make zip             # đóng gói bài nộp
make full            # cả quy trình: check, bench, figures, report, zip
\end{lstlisting}
% TODO(report): thời gian chạy thật của `make bench` trên từng máy (campaign.log); ghi chú riêng
% cho từng nền tảng nếu có (ví dụ cách cài matplotlib, compiler đã dùng).
% Phụ lục: ghi chú build. Nguồn: results/{apple-m5,LAP1184}/{campaign.log, machine.txt}, Makefile, mk/zz_pack.mk.

Trên \PCname{apple-m5}, đợt đo có 12\,152 lần chạy hoàn thành. Thời lượng 1 giờ 25 phút 45 giây gồm cả
2 phút hiệu chỉnh (643 lần chạy thử, không tính vào số trên). Lệnh đã dùng là
\texttt{python3 scripts/run\_bench.py --relax-yield --host apple-m5 --wait 30}, tương đương
\texttt{make bench} cộng lượt tùy chọn \texttt{make bench-relax}. \texttt{make bench-dry} in lưới và thời
gian ước tính mà không đo.
% nguồn: campaign.log (START argv=..., END duration=1:25:45); results/apple-m5/campaign_stdout.log ("calibrating 643
% iteration counts", "calibration done in 2:04")

Hai thí nghiệm phụ trên \PCname{apple-m5} (khoảng trống khi trao lock và đo lại lúc tải thấp) dùng
\texttt{experiments/handoff\_gap.cpp} và \texttt{experiments/run\_extra.py}. Chương trình
\texttt{handoff\_gap} chỉ cần các header trong \texttt{include/}, build bằng \texttt{mkdir -p
build/experiments \&\& c++ -Iinclude -std=c++20 -O2 -pthread experiments/handoff\_gap.cpp -o
build/experiments/handoff\_gap}. Lượt đo lại dùng \texttt{build/bench\_lock}, nên cần chạy
\texttt{make} trước. Lệnh \texttt{python3 experiments/run\_extra.py apple-m5} ghi
\texttt{handoff.csv}, \texttt{lowload.csv} và \texttt{run.log} vào \texttt{results/apple-m5/extra/}.
Số trong báo cáo tính bằng \texttt{scripts/report\_numbers/extra\_analyze.py}.
% nguồn: experiments/handoff_gap.cpp (include: bench_common.hpp, lock.h, platform.hpp, workload.hpp;
% omp.h chỉ khi _OPENMP); experiments/run_extra.py (docstring: Part A, Part B, Output)

Nền tảng gồm macOS 27.0.1, Apple clang 21.0.0, Python 3.14.8 (chạy đo chỉ cần thư viện chuẩn).
Apple clang không nhận \texttt{-fopenmp} nên không có baseline OpenMP. \texttt{matplotlib} và
\texttt{numpy} cài trong một môi trường ảo (\texttt{.venv}).
% nguồn: results/apple-m5/machine.txt

Trên \PCname{LAP1184}, nhóm chạy toàn bộ quy trình bằng \texttt{make full PYTHON=.venv/bin/python}
trong Ubuntu 20.04 trên WSL2. Đợt đo (\texttt{run\_bench.py --wait 10}) dài 1 giờ 20 phút 36 giây
(14\,322 lần chạy hoàn thành). Ubuntu 20.04 cần thêm g++ 11 từ PPA \texttt{ubuntu-toolchain-r/test} (g++ 9 mặc
định không có \texttt{<barrier>}). Nó cũng cần CMake 3.30 từ bản cài sẵn của Kitware (bản 3.16 của Ubuntu dưới mức
3.20 mà \texttt{CMakeLists.txt} cần), \texttt{python3-venv}, XeLaTeX và \texttt{latexmk}. Python 3.8.10
đủ cho các script.
\begin{itemize}[nosep]
  \item Git trên Windows mặc định đổi kết thúc dòng thành CRLF khi checkout, và \texttt{sh} trong WSL
    không chạy được \texttt{scripts/machine\_info.sh} khi đó. Runner bắt lỗi và ghi một
    \texttt{machine.txt} rỗng. Kho mã có \texttt{.gitattributes} giữ LF cho các file \texttt{.sh}.
    \texttt{machine.txt} của \PCname{LAP1184} được sinh lại bằng đúng lệnh của runner, 3 phút sau khi đợt
    đo kết thúc. Vì vậy dòng nguồn điện và load average trong đó là của lúc sinh lại.
  \item ThreadSanitizer báo data race giả trên \texttt{omp\_lock\_t}, vì libgomp không được instrument.
    Bản TSan của \texttt{test\_locks} và \texttt{test\_counter} bỏ qua lock OpenMP và in ra việc bỏ
    qua, như \texttt{test\_barrier} làm với \texttt{omp\_barrier}. Tham số \texttt{tsan\_uninstrumented}
    vẫn chạy nó để thấy cảnh báo.
  \item Không chạy \texttt{make bench-relax}, vì lựa chọn \texttt{isb}/\texttt{yield} chỉ có trên AArch64.
\end{itemize}
% nguồn: results/LAP1184/campaign.log (START argv=--wait 10; END duration=1:20:36, 01:41:24);
% results/LAP1184/machine.txt (date: 01:44:04); docs/part1.md mục 7.2.
Trên \PCname{C4M7Y670NT}, đợt đo (\texttt{run\_bench.py --wait 10}) dài 1 giờ 53 phút 10 giây (13\,244 lần
chạy hoàn thành), với macOS 26.7.1, Apple clang 21.0.0 và Python 3.9.6. Cách cài giống \PCname{apple-m5}.
% nguồn: results/C4M7Y670NT/campaign.log (START argv=--wait 10; END duration=1:53:10 ok=13244);
% results/C4M7Y670NT/machine.txt (python: Python 3.9.6).

% Bảng tên cài đặt. Nguồn: include/lock.h (PC_FOR_EACH_LOCK), bench/bench_barrier.cpp (danh sách barrier),
% bench/bench_counter.cpp (PC_COUNTERS_ATOMIC, PC_COUNTERS_LOCK), include/workload.hpp (pc::opaque).
\section{Tên các cài đặt}

Bảng~\ref{tab:impl:names} cho tên của mọi cài đặt, đúng như trong hình, bảng và log. Lock và barrier
thư viện chỉ dùng để so sánh.

\begin{table}[H]\centering\scriptsize\setlength{\tabcolsep}{3pt}
\begin{tabularx}{\linewidth}{@{}l>{\raggedright\arraybackslash}Xl@{}}
\toprule
\textbf{Tên} & \textbf{Mô tả} & \textbf{Hệ điều hành} \\
\midrule
\multicolumn{3}{@{}l}{\emph{Lock}} \\
\PCname{tas} & TAS: lặp \texttt{exchange(acquire)}, nhả bằng \texttt{store(release)} & mọi máy \\
\PCname{ttas} & TTAS: chờ bằng \texttt{load(relaxed)}, rồi mới \texttt{exchange} & mọi máy \\
\PCname{mcs} & MCS: hàng đợi FIFO, mỗi thread chờ trên nút riêng & mọi máy \\
\PCname{tas_bo_{sleep,spin,yield}} & TAS, gọi chính sách chờ sau mỗi lần thử hỏng (\texttt{sleep\_for} nhân đôi; quay ngẫu nhiên; quay 64 lần rồi \texttt{yield}) & mọi máy \\
\PCname{ttas_bo_{sleep,spin,yield}} & TTAS, gọi chính sách chờ khi thua \texttt{exchange} & mọi máy \\
\PCname{std_mutex}, \PCname{pthread_mutex} & \texttt{std::mutex}, \texttt{pthread\_mutex\_t} & mọi máy \\
\PCname{os_unfair_lock} & \texttt{os\_unfair\_lock} của macOS & macOS \\
\PCname{pthread_spinlock} & \texttt{pthread\_spinlock\_t} & Linux \\
\PCname{omp_lock} & \texttt{omp\_lock\_t} (cần bản dựng có OpenMP) & Linux \\
\midrule
\multicolumn{3}{@{}l}{\emph{Barrier}} \\
\PCname{sr} & sense-reversing barrier, chờ bằng quay & mọi máy \\
\PCname{sr_bo_{sleep,spin,yield}} & \PCname{sr}, gọi chính sách chờ như lock giữa hai lần đọc cờ & mọi máy \\
\PCname{sr_block} & \PCname{sr}, chờ chặn bằng \texttt{std::atomic::wait} và \texttt{notify\_all} & mọi máy \\
\PCname{std_barrier} & \texttt{std::barrier} (C++20) & mọi máy \\
\PCname{condvar} & mutex, condition variable và bộ đếm thế hệ & mọi máy \\
\PCname{pthread_barrier} & \texttt{pthread\_barrier\_t} (macOS không có) & Linux \\
\PCname{omp_barrier} & \texttt{\#pragma omp barrier} (cần bản dựng có OpenMP) & Linux \\
\midrule
\multicolumn{3}{@{}l}{\emph{Counter}} \\
\PCname{faa}, \PCname{faa_seqcst} & \texttt{fetch\_add}, thứ tự \texttt{relaxed} hoặc \texttt{seq\_cst} & mọi máy \\
\PCname{cas}, \PCname{cas_seqcst} & vòng CAS không backoff, thứ tự \texttt{relaxed} hoặc \texttt{seq\_cst} & mọi máy \\
\PCname{cas_{sleep,spin,yield}} & vòng CAS, gọi chính sách chờ sau mỗi lần CAS hỏng & mọi máy \\
\PCname{lock_X} & biến \texttt{long} trong critical section của lock \PCname{X} (mọi lock ở trên) & như lock \PCname{X} \\
\bottomrule
\end{tabularx}
\caption{Tên các cài đặt. Hai máy macOS (\PCname{apple-m5}, \PCname{C4M7Y670NT}) có 12 lock, 7 barrier
và 19 counter. \PCname{LAP1184} (Linux, có OpenMP) có 13 lock, 9 barrier và 20 counter.}
\label{tab:impl:names}
\end{table}
% nguồn: results/{apple-m5,C4M7Y670NT}/test.log (12 lock, 7 barrier, 19 counter); results/LAP1184/test.log (13, 9, 20).

Benchmark và bộ kiểm tra cho các giá trị mà compiler có thể gộp hoặc xoá đi qua \texttt{pc::opaque()}:
đối số $N$ của mỗi lời gọi \texttt{is\_prime}, địa chỉ counter ở mỗi vòng của bài kiểm tra counter, và
giá trị đọc ra của biến đếm trong bài kiểm tra lock. Hàm này chỉ là một lệnh \texttt{asm volatile} rỗng
với ràng buộc \texttt{+r}. Compiler phải coi giá trị trong thanh ghi đã bị sửa, nên không chứng minh được
hai lời gọi giống nhau và không gộp hay xoá chúng. Hàm không sinh lệnh máy nào.
% nguồn: include/workload.hpp (opaque); bench/bench_lock.cpp:81, bench/bench_counter.cpp:102,
% tests/test_counter.cpp:326, tests/test_locks.cpp:116.

\section{Kết quả kiểm tra, đợt đo và hình bổ sung}
% Bảng tóm tắt kết quả kiểm tra tính đúng đắn (phụ lục), trỏ tới từ corr_strategy.tex.
\begin{table}[H]\centering\scriptsize\setlength{\tabcolsep}{4pt}
\begin{tabularx}{\linewidth}{@{}>{\raggedright\arraybackslash}Xlll@{}}
\toprule
 & \textbf{\PCname{apple-m5}} & \textbf{\PCname{C4M7Y670NT}} & \textbf{\PCname{LAP1184}} \\
\midrule
Số thread của stress test & 1, 2, 4, 10, 20 & 1, 2, 4, 12, 24 & 1, 2, 4, 12, 24 \\
\texttt{make test}: lock, barrier, counter đạt & 12/12, 7/7, 19/19 & 12/12, 7/7, 19/19 & 13/13, 9/9, 20/20 \\
\quad barrier với schedule fuzzing đạt & 5/5 & 5/5 & 5/5 \\
\quad episode barrier được kiểm & 2\,198\,144 & 2\,198\,201 & 2\,085\,088 \\
\quad lần tăng counter được kiểm & 121\,573\,312 & 144\,627\,911 & 148\,879\,205 \\
\texttt{make tsan}: lock, barrier, counter đạt & 12/12, 7/7, 19/19 & 12/12, 7/7, 19/19 & 12/12, 8/8, 19/19 \\
\quad bị bỏ qua & handoff của \PCname{std_barrier} & handoff của \PCname{std_barrier} & \PCname{omp_lock}, \PCname{omp_barrier}, \PCname{lock_omp_lock} \\
\quad cảnh báo data race & 0 & 0 & 0 \\
\texttt{make mutation}: biến thể bị bắt & 14/14 & 14/14 & 14/14 \\
\bottomrule
\end{tabularx}
\caption{Kết quả kiểm tra tính đúng đắn trên ba máy, mỗi máy một lượt. Episode barrier chỉ tính bản
thường. Lý do bỏ qua trong bản TSan nằm ở các Mục~\ref{sec:corr:lock}--\ref{sec:corr:counter}.}
\label{tab:corr:summary}
\end{table}
% nguồn: results/{apple-m5,C4M7Y670NT,LAP1184}/test.log (dòng PASS; tổng episode và lần tăng cộng bằng awk;
% episode barrier chỉ cộng các dòng PASS của phần "barrier correctness test (plain)"),
% tsan.log (dòng SKIP, note), mutation.log.

\CampaignTables{apple-m5}
\PCtable{apple-m5}{locksingle}
\PCtable{apple-m5}{sleeplatency}
\PCfigure{apple-m5}{lock_cs_length}
\clearpage
\section{Bảng số liệu của hai máy còn lại}
Các bảng dưới đây cùng dạng với bảng của \PCname{apple-m5} trong mục Phương pháp đo và mục Kết quả và
phân tích; mục So sánh giữa các máy trích số từ chúng.
\newcommand{\OtherHostTables}[1]{\PCtable{#1}{campaign}\PCtable{#1}{lockunit}\PCtable{#1}{locktime}\PCtable{#1}{barrier}\PCtable{#1}{countertime}}
\ForEachOtherHost{\OtherHostTables}
\section{Cấu trúc mã nguồn}
% TODO(report): kiểm lại bảng này với nội dung thật của file zip trước khi nộp.
\begin{table}[H]\centering\small
\begin{tabular}{@{}lp{10.2cm}@{}}
\toprule
\textbf{Đường dẫn} & \textbf{Nội dung} \\
\midrule
\texttt{include/lock.h} & các lock của phần 1 (TAS, TTAS, biến thể backoff, MCS) và lock thư viện \\
\texttt{include/barrier.hpp} & sense-reversing barrier và các barrier dùng để so sánh \\
\texttt{include/counter.hpp} & các counter của phần 3 \\
\texttt{include/backoff.hpp} & các chính sách chờ dùng chung cho cả ba phần \\
\texttt{include/platform.hpp} & cache line, lệnh chờ trong vòng spin, đồng hồ, thông tin máy \\
\texttt{include/workload.hpp} & \texttt{is\_prime} và hàm \texttt{opaque} chống trình tối ưu \\
\texttt{include/bench\_common.hpp} & đọc tham số, cổng xuất phát, định dạng dòng CSV \\
\texttt{tests/} & bài kiểm tra tính đúng đắn của ba phần \\
\texttt{bench/} & ba chương trình đo và \texttt{sleep\_latency} \\
\texttt{mk/} & các mục tiêu \texttt{make} của từng phần và của việc đóng gói \\
\texttt{scripts/run\_bench.py} & chạy toàn bộ đợt đo, ghi số liệu thô vào \texttt{results/} \\
\texttt{scripts/plot.py} & vẽ mọi hình, sinh mọi bảng của báo cáo từ \texttt{results/} \\
\texttt{scripts/machine\_info.sh} & ghi lại cấu hình máy \\
\texttt{scripts/fake\_bench.py} & chương trình giả để thử công cụ, không tạo ra số đo \\
\texttt{results/<máy>/} & số liệu thô của từng máy (CSV), cấu hình máy, nhật ký đợt đo \\
\texttt{report/} & mã nguồn báo cáo, hình và bảng sinh tự động \\
\bottomrule
\end{tabular}
\caption{Cấu trúc thư mục của bài nộp.}\label{tab:layout}
\end{table}

\end{document}
